%%% CV de Julie Le Gallo %%%

%%% Pr\'eambule %%%

\documentclass[a4paper,12pt]{article}

\usepackage[french]{babel}	% faire du fran\'eais
\usepackage[T1]{fontenc}		% c\'esure
\usepackage[latin1]{inputenc}	% utilisation des touches accentu\'ees
\usepackage{pifont}				% inclure des images dingbats
\usepackage{url}				% citer des URL
\usepackage[left = 2.5cm, right = 2.5cm, top = 2.5cm, bottom = 2.5cm]{geometry}	% changer la mise en page
\usepackage{soul}				% pour pouvoir souligner plusieurs lignes
\usepackage{longtable}			% pour des tableaux sur plusieurs pages
\usepackage{enumitem}

\AtBeginDocument{\renewcommand{\labelitemi}{\textbullet}}	    % red\'efinition des listes
\AtBeginDocument{\renewcommand{\arraystretch}{1.1}}		    % ajout d'espace entre les lignes des tableaux

%%% D\'efinition des commande pour les titres %%%

\newcommand{\bigtitre}[1]{%			% titre g\'en\'eral
	\begin{center}
		\textbf{\LARGE \textsc{#1}}
	\end{center}
	\medskip
	}
	
\newcommand{\titre}[1]{%			% titre
	\bigskip
	\noindent \rule[2mm]{\textwidth}{0.5pt}
	\textbf{\textsc{\large #1}} \\
	\noindent \rule[1mm]{\textwidth}{0.5pt}
	}
	
\newcommand{\smalltitre}[1]{%		% sous-titre
	\noindent \textbullet~\textsc{\textbf{#1}}
	\smallskip
	}
	
%%% d\'ebut du document %%%
 
\begin{document}

%%% informations g\'en\'erales

\begin{center}
	\textbf{{\Large CURRICULUM VITAE}} \smallskip
	\\ \textbf{{\large Avril 2026}} 	\bigskip \bigskip
	\\ \textbf{{\Large Julie LE GALLO}}
\end{center}

\bigskip

%%% Fonction et \'etablissement actuels %%%

\titre{Fonction et \'etablissement actuels}

\begin{itemize}[leftmargin=*]
	\item Professeure de l'enseignement sup\'erieur agricole, l'Institut Agro Dijon
	\item Membre du \textit{Centre d'Economie et de Sociologie appliqu\'ees  \`a l'Agriculture et aux Espaces Ruraux} (CESAER, UMR 1041)
	\item Membre associ\'ee : \\
	\ding{219} \textit{Grupo de Econom\'ia Regional y Espacial} (ECONRES, Universidad Aut\'onoma de Madrid, Espagne)
\end{itemize} 

\bigskip
 
\textit{Adresse professionnelle}
\par \quad{CESAER}
\par \quad{26, bd Docteur Petitjean}
\par \quad{BP 87999}
\par \quad{21079 Dijon Cedex, France} \smallskip
\par \textit{T\'el\'ephone}: + 33 3 80 77 23 66
\par \textit{Courriel}: \url{julie.le-gallo@institut-agro.fr} 
\par \textit{Page personnelle}: \url{https://sites.google.com/site/legallopage/} \\

%%% Exp\'erience professionnelle %%%

\titre{Exp\'erience professionnelle}

% mise en forme des tableaux 
% \annee est la largeur de la 1ere colonne contenant l'ann\'e. Elle est ici d\'efinie comme \'etant la largeur du texte " 2003 - 2006 ".

\newlength{\annee}
\settowidth{\annee}{2003 - 2006}

% \texte est la largeur de la deuxi\'eme colonne. Elle est d\'efinie comme \'etant la largeur de la page moins celle de la premi\'ere colonne.
% 2\tabcolsep est la largeur de l'espacement entre les colonnes.

\newlength{\texte}
\setlength{\texte}{\textwidth} \addtolength{\texte}{-\annee} 
	\addtolength{\texte}{-1\tabcolsep}

\noindent \begin{longtable}{@{}p{\annee}p{\texte}@{}}
2022 - 		& Professeure de l'enseignement sup\'erieur agricole \`a l'Institut Agro Dijon \\
2015 - 2021	& Professeure de l'enseignement sup\'erieur agricole \`a Agrosup Dijon \\
2008 - 2015	& Professeure des Universit\'es \`a l'U.F.R. Sciences Juridiques, Economiques, Politiques et de Gestion, Universit\'e de Franche-Comt\'e \\
2006 - 2008	& Professeure des Universit\'es \`a l'I.U.T. de Belfort-Montb\'eliard, D\'epartement Techniques de Commercialisation, Universit\'e de Franche-Comt\'e \\
2003 - 2006	& Ma\^itresse de Conf\'erences \`a l'Universit\'e Montesquieu-Bordeaux IV \\
2002 - 2003	& Chercheuse invit\'ee au \textit{Regional Economics Applications Laboratory}, Universit\'e d'Illinois en Urbana-Champaign (Etats-Unis) \\
2001 - 2002	& A.T.E.R. \`a l'U.F.R. de Sciences Economiques et de Gestion, Universit\'e de Bourgogne \\
1998 - 2001	& Allocataire-monitrice \`a l'U.F.R. de Sciences Economiques et de Gestion, Universit\'e de Bourgogne
\end{longtable}

%%% Titres universitaires %%%

\titre{Titres universitaires}

\noindent - \textbf{Premier concours d'agr\'egation de l'Enseignement Sup\'erieur en Sciences Economiques}, 2006, \'equivalence Habilitation \`a Diriger des Recherches \\
- \textbf{Doctorat en Analyse et Politique Economiques} \newline \ding{219} \og{}\textit{Disparit\'es g\'eographiques et convergence des r\'egions europ\'eennes : une approche par l'\'econom\'etrie spatiale}\fg{}, th\`ese soutenue le 21 Mai 2002 \`a l'Universit\'e de Bourgogne. \textbf{Mention tr\`es honorable, f\'elicitations du jury \`a l'unanimit\'e et proposition pour prix de th\`ese} \\ 
%\textit{Jury} & C.~Baumont (directrice de th\`ese), G.~Benhayoun (pr\'esident), C.~Ertur (suffrageant), B.~Fingleton (suffrageant), R.J.G.M.~Florax (rapporteur), P.~Morin (rapporteur), M.-C.~Pichery (co-directrice de th\`ese) \\ 
%1998 & \textbf{DEA en Analyse et Politique Economiques (mention Economie Math\'ematique et Econom\'etrie)} \newline \ding{219} \og{}\textit{L'espace dans les mod\`eles \'econom\'etriques}\fg{}, soutenu en septembre 1998 \`a l'Universit\'e de Bourgogne. %Mention Tr\`es Bien (Premier prix) \\ 

%%% Prix et distinctions scientifiques %%%

\titre{Prix et distinctions scientifiques}

\noindent - \textbf{Musgrave Prize 2025} for the best paper published in \textit{National Tax Journal},  2025 \\
- \textbf{Prix \og{}Excellence in Reviewing\fg{} pour la revue \textit{Geographical Analysis}}, 2014 \\
- \textbf{Fellow} \'elue de la \textbf{\textit{Spatial Econometrics Association}}, 2010 \\
- \textbf{Prix \og{}Best referee\fg{} pour la revue \textit{Regional Studies}}, 2006 \\
- \textbf{2\ieme{} au prix EPAINOS}, Prix jeune chercheur d\'ecern\'e par l'\textit{European Regional Science Association} (ERSA), 2005 \\%pour le meilleur papier pr\'esent\'e au Congr\`{e}s annuel de l'ERSA. Papier prim\'e : \og{}Spatial and sectoral productivity convergence between European regions, 1975-2000\fg{} (avec S.~Dall'erba), Vrije Universiteit, Amsterdam (Pays-Bas), 23-27 ao\^ut \\
- \textbf{2\ieme{} au prix EPAINOS}, Prix jeune chercheur d\'ecern\'e par l'\textit{European Regional Science Association} (ERSA), 2004 \\% pour le meilleur papier pr\'esent\'e au Congr\`{e}s annuel de l'ERSA. Papier prim\'e : \og{}The evolution of the spatial and sectoral patterns in Ile-de-France over 1978-1997\fg{} (avec R.~Guillain et C.~Boiteux-Orain), Porto (Portugal), 25-29 ao\^ut \\
- \textbf{Prix Philippe Aydalot}, Prix \og{}Jeune Chercheur en Science R\'egionale\fg{} d\'ecern\'e par l'\textit{Association de Science R\'egionale De Langue Fran\c{c}aise}, 2003 \\
- \textbf{Prix Vouters}, Prix de th\`ese d\'ecern\'e par l'U.F.R. de Sciences Economiques et de Gestion de l'Universit\'e de Bourgogne, 2003

%%% S\'ejours \'e l'\'etranger %%%

\titre{S\'ejours \`a l'\'etranger}

\noindent - \textbf{S\'ejour de recherche} d'une semaine \`a l'Universidad Auton\'oma de Madrid (Espagne), invit\'ee par le Professeur C.~Chasco, 2009 \\
- \textbf{S\'ejour de recherche} d'un mois au SAL (S\textit{patial Analysis Laboratory}) \`a l'Universit\'e d'Illinois en Urbana-Champaign (Etats-Unis), invit\'ee par le Professeur L.~Anselin, 2004 \\
- \textbf{S\'ejour post-doctoral} en tant que chercheur invit\'e au REAL (\textit{Regional Economics Applications Laboratory}) \`a l'Universit\'e d'Illinois en Urbana-Champaign (Etats-Unis), invit\'ee par le Professeur L.~Anselin, 2002-2003 \\
- \textbf{S\'ejour pr\'e-doctoral} de deux mois, CORE et D\'epartement de G\'eographie, Universit\'e de Louvain-la-Neuve (Belgique), invit\'ee par le Professeur D.~Peeters, 2000

\newpage

%%% Responsabilit\'es collectives %%%

\bigtitre{Responsabilit\'es collectives}

%%% Responsabilit\'es d'expertise et d'\'evaluation scientifique %%%

\titre{Responsabilit\'es d'expertise et d'\'evaluation scientifique}

\smalltitre{Au niveau national}

\noindent \begin{longtable}{@{}p{\annee}p{\texte}@{}}
2018 - 		& Evaluation de projets de th\`ese CIFRE pour l'ANRT \newline
\ding{219} 2018, 2022 \\
2006 - 		& Evaluation de projets de recherche pour l'\textbf{Agence Nationale pour la Recherche (ANR)} \newline
\ding{219} 2006 (1), 2010 (1), 2022 (2), 2023 (1), 2024 (1 projet + membre du comit\'e CE26 - Individus, entreprises, march\'es, finance, management), 2025 (membre du comit\'e CE26 - Individus, entreprises, march\'es, finance, management) \\
2015 - 2025	& Expert pour le \textbf{Haut Conseil de l'Evaluation de la Recherche et de l'Enseignement Sup\'erieur (HCERES)} \newline
\ding{219} 2020, 2021, 2024, 2025 - Pr\'esidence d'un jury d'\'evaluation d'une unit\'e de recherche \newline
\ding{219} 2015, 2019, 2024 - Jury d'\'evaluation d'une unit\'e de recherche \newline
\ding{219} 2017 - Evaluation de mentions de master \\
2023 - 2025	& Membre du Comit\'e Scientifique en charge de l'\'evaluation de la seconde phase de l'Exp\'erimentation Territoires Z\'ero Ch\^omeurs de Longue Dur\'ee \\
2019   		& Evaluation d'un projet de recherche pour l'Institut d'Etudes Avanc\'ees de Paris \\
2013   		& Membre du jury d'experts pour l'attribution de la Prime d'Excellence Scientifique en section 5 (Economie) \\
2010  		& Evaluation d'un projet de recherche pour l'\textbf{Universit\'e du Littoral C\^ote d'Opale} \\
2008 - 2014	& Expert r\'egulier pour l'\textbf{Agence d'Evaluation de la Recherche et de l'Enseignement Sup\'erieur (AERES)} \newline 
\ding{219} 2008, 2009 - Evaluation d'unit\'es de recherche \newline 
\ding{219} 2008, 2009, 2010, 2012, 2013, 2014 - Evaluation de formations (licence et master) dispens\'ees dans les universit\'es \\
2008   		& Evaluation d'un projet de recherche pour : \newline
\ding{219} la \textbf{r\'egion Provence-Alpes-C\^ote d'Azur}, Direction de l'Economie R\'egionale, de l'Innovation et de l'Enseignement Sup\'erieur \newline
\ding{219} la \textbf{facult\'e Jean Monnet} de l'Universit\'e Paris-Sud 11 \\
2007 - 2010	& Membre nomm\'e de la 5\ieme{} section du \textbf{Conseil National des Universit\'es (CNU)}, coll\`{e}ge A \\
2007   		& Expert pour la \textbf{Mission Scientifique, Technique et P\'edagogique (MSTP)}, DSPT 7, D\'epartement des Sciences de la Soci\'et\'e, Economie \\
\end{longtable}

\newpage

\smalltitre{Au niveau international}

\noindent \begin{longtable}{@{}p{\annee}p{\texte}@{}}
2025   	& Pr\'esidence du Comit\'e \textit{Fundac\~ao para a Ci\^encia e a Tecnologia} (Portugal) pour la section\textit{ Economics \& Business} \\
2024   	& Pr\'esidence du Comit\'e \textit{Fundac\~ao para a Ci\^encia e a Tecnologia} (Portugal) pour la section\textit{ Economics \& Business} \\
2024   	& Evaluation d'un projet de recherche pour ORA8, Open Research Area for the Social Sciences \\
2024   	& Evaluation de deux projets de recherche pour EUTOPIA Science and Innovation Fellowship Program  \\
2021   	& Evaluation d'un projet de recherche pour EUTOPIA Science and Innovation Fellowship Program \\
2019  	& Expert pour l'\textbf{European Research Council (ERC)} \newline \ding{219} Evaluation d'un 'ERC Starting Grant' \\
2015 	& Evaluation d'un projet de recherche pour la \textbf{Czech Science Foundation} \\
		& Evaluation d'un projet de recherche pour le \textbf{Conseil de recherches en sciences humaines du Canada} \\
2014  	& Expert pour l'\textbf{Economic and Social Research Council (ESRC)}\\
2014  	& Expert pour l'\textbf{European Research Council (ERC)} \newline \ding{219} Evaluation d'un 'ERC Starting Grant' \\
2013  	& Expert pour l'\textbf{European Science Foundation (ESF)} \\
2012  	& Expert pour l'\textbf{European Science Foundation (ESF)} \\
2010  	& Expert pour l'\textbf{European Science Foundation (ESF)} \\
2009 - 2010 & Membre du \textbf{jury d\'ecernant le prix EPAINOS} \newline \ding{219} Prix jeune chercheur d\'ecern\'e chaque ann\'ee par l'\textit{European Regional Science Association} (ERSA) \\
2009 	& Evaluation d'un projet de recherche pour le \textbf{Conseil de recherches en sciences humaines du Canada} \\
2008 	& Evaluation d'un projet de recherche pour le \textbf{Austrian Science Fund}, Department for Humanities and Social Sciences \\
2007 	& Evaluation d'un projet de recherche pour le \textbf{Conseil de recherches en sciences humaines du Canada} \\
2004   	& Evaluation d'un projet de recherche pour l'\textbf{Universit\'e du Luxembourg}
\end{longtable}

%%% Responsabilit\'es \'editoriales %%%

\titre{Responsabilit\'es \'editoriales}

%%% Co-\'editeur %%%

\smalltitre{Editeur}

\noindent \begin{longtable}{@{}p{\annee}p{\texte}@{}}
2026 - & Editeur, \textit{Journal of Regional Science} \\
2025 - & Editeur en chef, \textit{Journal of Geographical Systems} \\
2024 - & Editeur de section, \textit{Frontiers in Sustainable Cities - section Urban Economics} \\
2008 - & Editeur associ\'e, \textit{R\'egion et D\'eveloppement} \\
2019 - 2024 & Co-\'editeur, \textit{Papers in Regional Science} \\
2007 - 2021 & Co-\'editeur, \textit{Spatial Economic Analysis} \\
2010 - 2015 & Co-\'editeur, \textit{Revue d'Economie R\'egionale et Urbaine} \\
\end{longtable}

%%% Comit\'e \'editorial %%%

\smalltitre{Comit\'es \'editoriaux}
	
\noindent \begin{longtable}{@{}p{\annee}p{\texte}@{}}
2024 - 2025 & Membre du comit\'e \'editorial, \textit{Journal of Regional Science} \\
2024 - & Membre du comit\'e \'editorial, \textit{Econometrics} \\
2023 - & Membre du comit\'e \'editorial, \textit{Journal of Korean Planning Association} \\
2019 - 2024 & Membre du comit\'e \'editorial, \textit{Journal of Spatial Econometrics} \\
2019 - & Membre du comit\'e \'editorial, \textit{Revue Economique} \\
2019 - & Membre du comit\'e \'editorial, \textit{Investigaciones Regionales - Journal of Regional Research} \\
2019 - & Membre du comit\'e \'editorial, \textit{Journal of Geographical Systems} \\
2016 - 2019 & Membre du comit\'e \'editorial, \textit{Geopolitics under Globalization} \\
2016 - & Membre du comit\'e \'editorial, \textit{International Journal of Urban Sciences} \\
2015 - & Membre du comit\'e \'editorial, \textit{Regional Studies, Regional Science} \\
2014 - & Membre du comit\'e \'editorial, \textit{Economics and Business Letters} \\
2013 - 2014 & Membre du comit\'e \'editorial, \textit{ISRN Economics} \\
2010 - & Membre du comit\'e \'editorial, \textit{International Regional Science Review} \\
2009 - 2015 & Membre du comit\'e \'editorial, \textit{Papers in Regional Science} \\
2007 - & Membre du comit\'e \'editorial, \textit{Review of Regional Studies} \\
2007 - & Membre du comit\'e \'editorial, \textit{Letters in Spatial and Resource Science}
\end{longtable}

%%% Arbitre %%%

\smalltitre{Arbitre} \\

\noindent \textit{Agricultural Economics} (1), \textit{Annals of Regional Science} (20), \textit{Applied Economics} (7), \textit{Applied Economics Letters} (3), \textit{Applied Geography} (1), \textit{Applied Spatial Analysis and Policy} (3), \textit{B.E. Journal of Economic Analysis and Policy} (1), \textit{British Food Journal} (1), \textit{Buildings} (1), \textit{Case Studies in Business, Industry and Government Statistics} (1), \textit{China Economic Review} (1), \textit{Computational Statistics and Data Analysis} (1), \textit{Computers, Environment and Urban Systems} (2), \textit{Economic Geography} (2), \textit{Economic Journal} (1), \textit{Economic Modelling} (7), \textit{Economic Trends and Economic Policy} (1), \textit{Economics Bulletin} (3), \textit{Economics and Business Letters} (2), \textit{Empirical Economics} (10), \textit{Environment and Development Economics} (1), \textit{Environment and Planning A/B/C} (14), \textit{European Economic Review} (3), \textit{European Journal of Development Research} (1), \textit{European Review of Agricultural Economics} (4), \textit{Food Policy} (1), \textit{Forestry} (1), \textit{Geografiska Annaler B} (1), \textit{Geographical Analysis} (15), \textit{Geography Journal} (1), \textit{Geospatial Health} (1), \textit{Growth and Change} (4), \textit{Heliyon} (1), \textit{International Journal of Geo-Information } (1), \textit{International Economics} (1), \textit{International Economic Review} (2), \textit{International Journal of Manpower} (1), \textit{International Regional Science Review} (19), \textit{International Tax and Public Finance} (1), \textit{International Trade Journal} (1), \textit{ISRN Economics} (3), \textit{Italian Journal of Regional Science} (1), \textit{Investigaciones Regionales} (1), \textit{Journal of Applied Econometrics} (1), \textit{Journal of Applied Economics} (1), \textit{Journal of Criminal Justice} (2), \textit{Journal of Economic Geography} (6), \textit{Journal of Economic Surveys} (1), \textit{Journal of Economic Inequality} (2), \textit{Journal of Economic Studies} (1), \textit{Journal of Environmental Economics and Management} (1), \textit{Journal of Environmental Management} (3), \textit{Journal of Forest Economics} (1), \textit{Journal of Geographical Systems} (18), \textit{Journal of Human Development} (1), \textit{Journal of Policy Analysis and Management} (1), \textit{Journal of Regional Science} (15), \textit{Journal of Spatial Econometrics} (2), \textit{Journal of Urban Economics} (3), \textit{Land} (2), \textit{Land Use Policy} (1), \textit{Letters in Spatial and Resource Science} (7), \textit{Mathematics} (1),\textit{Networks and Spatial Econometrics} (2),  \textit{Open Journal of Statistics} (1), \textit{Oxford Bulletin of Economics and Statistics} (1),\textit{Papers in Regional Science} (30), \textit{Plos One} (1), \textit{Public Finance Review} (1), \textit{Real Estate Economics} (1), \textit{Regional Science and Urban Economics} (7), \textit{Regional Science Policy and Practice} (1), \textit{Regional Studies} (22), \textit{Regional Studies, Regional Science} (8), \textit{Research in Transportation Business and Management} (1), \textit{Research Policy} (5), \textit{Resource and Energy Economics} (1), \textit{Review of Agricultural, Food and Environmental Studies} (1), \textit{Review of Income and Wealth} (1), \textit{Review of Regional Studies} (24), \textit{Scottish Journal of Political Economy} (1), \textit{Social Indicators Research} (1), \textit{Socio-Economic Planning Sciences} (1), \textit{Socio-Economic Review} (1), \textit{Spanish Economic Review} (1), \textit{Spatial Economic Analysis} (10), \textit{Spatial Statistics} (11), \textit{Sustainability} (1), \textit{Technological Forecasting and Social Change} (1), \textit{The Econometrics Journal} (1), \textit{The Social Science Journal} (1), \textit{The World Economy} (1), \textit{Tijdschrift voor Economische en Sociale Geografie} (4), \textit{Transportation Research Record} (1), \textit{Urban Studies} (5), \textit{Urban Studies Research}, \textit{World Development} (2) \medskip 

\noindent \textit{Annales d'Economie et de Statistique} (3), \textit{Cahiers d'Economie de l'Innovation} (1), \textit{Cahiers d'Economie et de Sociologie Rurales} (2), \textit{Cahiers Economiques de Bruxelles} (1), \textit{Cahiers Scientifiques du Transport} (2), \textit{Cyberg\'eo} (1), \textit{Economie et Pr\'evision} (6), \textit{Economie et Statistique} (3), \textit{Economie Internationale} (2), \textit{Economie Publique} (1), \textit{Economie Rurale} (2), \textit{Management International} (2), \textit{Recherches Economiques de Louvain} (2), \textit{R\'egion et D\'eveloppement} (11), \textit{Revue d'Economie R\'egionale et Urbaine} (2), \textit{Revue de l'OFCE} (1), \textit{Revue d'Economie Politique} (2), \textit{Revue Economique} (18), \textit{Revue Internationale des Francophonies} (1), \textit{Travail et Emploi} (1)

%%% Responsabilit\'es p\'edagogiques %%%

\titre{Responsabilit\'es p\'edagogiques}

\noindent \begin{longtable}{@{}p{\annee}p{\texte}@{}}
2024 - 		& Responsable des \textbf{stages \`a l'international de deuxi\`{e}me ann\'ee de formation par apprentissage d'ing\'enieurs, sp\'ecialit\'e agronomie}, Institut Agro Dijon\newline
\ding{219} \textit{Missions p\'edagogiques} : accompagnement \`a la recherche de stages, validation p\'edagogique, suivi des stages des \'etudiants, \'evaluation \medskip \\ 
2019 - 2023	& Co-Responsable du master international "\textbf{Data Analyst in Spatial and Environmental Economics}", Universit\'e de Bourgogne Franche-Comt\'e \newline
\ding{219} \textit{Missions p\'edagogiques} : organisation et articulation des cours, organisation des interventions des intervenants ext\'erieurs, accompagnement \`a la recherche de stages, suivi des stages des \'etudiants, suivi des m\'emoires de recherche, promotion de la formation \medskip \\ 
2016 - 2021	& Responsable des \textbf{stages \`a l'international de deuxi\`{e}me ann\'ee de formation initiale d'ing\'enieurs, sp\'ecialit\'e agronomie}, AgroSup Dijon\newline
\ding{219} \textit{Missions p\'edagogiques} : accompagnement \`a la recherche de stages, validation p\'edagogique, suivi des stages des \'etudiants \medskip \\ 
2016 - 2019 & Responsable de l'unit\'e p\'edagogique \og{}Politiques publiques\fg{}, AgroSup Dijon \newline
\ding{219} \textit{Missions p\'edagogiques} : organisation et coordination des enseignements d'\'economie des politiques publiques \`a AgroSup Dijon \medskip \\ 
2012 - 2015	& Responsable de la sp\'ecialit\'e \og{}\textbf{Charg\'e d'Etudes Economiques}\fg{} du Master mention \og{}\textbf{Economie - Gestion}\fg{}, Universit\'e de Franche-Comt\'e \newline
\ding{219} \textit{Missions p\'edagogiques} : organisation et articulation des cours, organisation des interventions des professionnels, accompagnement \`a la recherche de stages, suivi des stages des \'etudiants, suivi des m\'emoires de recherche \medskip \\  
2008 - 2009 & Responsable de la sp\'ecialit\'e \og{}\textbf{Expertise Economique}\fg{} du Master mention \og{}\textbf{Economie de la Firme et des March\'es}\fg{}, Universit\'e de Franche-Comt\'e  \newline
\ding{219} \textit{Missions p\'edagogiques} : organisation et articulation des cours, organisation des interventions des professionnels, accompagnement \`a la recherche de stages, suivi des stages des \'etudiants, suivi des m\'emoires de recherche \\ 
\end{longtable}

%%% Responsabilit\'es scientifiques %%%

\titre{Autres responsabilit\'es}

\smalltitre{Animation de la recherche}

\noindent \begin{longtable}{@{}p{\annee}p{\texte}@{}}
2020 -  		& Co-responsable des s\'eminaires en ligne d'\'economie g\'eographique du Centre d'Economie et de Sociologie appliqu\'ees \`a l'Agriculture et aux Espaces Ruraux (CESAER, UMR1041, Institut Agro Dijon, INRAE), \url{https://sites.google.com/view/secgo/home} \\
2015 -  2020	& Co-responsable de l'axe \textbf{Dynamiques et am\'enagement des territoires} du Centre d'Economie et de Sociologie appliqu\'ees \`a l'Agriculture et aux Espaces Ruraux (CESAER, UMR1041, AgroSup Dijon, INRAE) \\
2006 - 2009 	& Responsable de l'axe \textbf{Econom\'etrie Appliqu\'ee} du Centre de Recherche sur les Strat\'egies Economiques (CRESE, EA 3190) \\
\end{longtable} 

\smalltitre{Responsabilit\'es \'electives} \\

\noindent \textit{\textbf{Au sein d'AgroSup Dijon, devenu Institut Agro Dijon en 2022, 2015 - }}

\noindent \begin{longtable}{@{}p{\annee}p{\texte}@{}}
	2022 - 		& Membre suppl\'eant \'elu du \textbf{Conseil des Etudes et de la Vie Etudiante} \\
	2021 - 2023	& Directrice du \textbf{D\'epartement des Sciences Humaines et Sociales} \\	
	2019 - 2021	& Membre suppl\'eant \'elu du \textbf{Conseil des Enseignants} \\
	2019 - 2021	& Membre suppl\'eant \'elu du \textbf{Conseil des Etudes et de la Vie Etudiante} \\
	2016 - 2021	& Membre suppl\'eant \'elu du \textbf{Conseil d'Administration}, r\'e\'elue en 2019 \\	
	2016 - 2019	& Membre titulaire \'elu du \textbf{Conseil des Enseignants} \\	
\end{longtable}

\noindent \textit{\textbf{Au sein de l'Universit\'e de Franche-Comt\'e, 2008 - 2015}}

\noindent \begin{longtable}{@{}p{\annee}p{\texte}@{}}
	2008 - 2014	& Membre \'elu du \textbf{Conseil d'Administration} de l'U.F.R. de Sciences Juridiques, Economiques, Politiques et de Gestion (r\'e\'elue en 2010) \\
	2013 		& Membre d'un \textbf{Comit\'e de S\'election} \\
	2008 - 2012 & Membre \'elu du \textbf{Conseil Scientifique} (coll\`{e}ge A, domaine Droit-Economie-Gestion) \\
	& Membre du \textbf{Conseil de la Documentation} \\
	2010 		& Membre de \textbf{Comit\'es de S\'election} \\
	2008 - 2009 & Membre du \textbf{Conseil de l'Ecole Doctorale Louis Pasteur} \\
\end{longtable}


\smalltitre{Jurys de concours} 

\noindent \begin{longtable}{@{}p{\annee}p{\texte}@{}}
2025	& Membre d'un \textbf{Comit\'e de S\'election}, Universit\'e Paris-Nanterre\\
2024	& Pr\'esidente d'un \textbf{Comit\'e de S\'election}, Universit\'e Autonome de Barcelone (Espagne) \\
		& Membre d'un \textbf{Comit\'e de S\'election}, Universit\'e Lyon 2\\
		& Membre d'un \textbf{Comit\'e de S\'election}, Universit\'e de Bordeaux \\
		& Membre d'un \textbf{Comit\'e de S\'election}, Universit\'e de Clermont-Ferrand \\
2023	& Membre d'un \textbf{Comit\'e de S\'election}, Universit\'e Jean Monnet  Saint-Etienne \\
		& Membre d'un \textbf{Comit\'e de S\'election}, Universit\'e de Rennes \\
		& Membre d'un \textbf{Comit\'e de S\'election}, Ecole Nationale des Travaux Publics de l'Etat \\
2022	& Membre d'un \textbf{Comit\'e de S\'election}, Institut de Hautes Etudes en Administration Publique, Universit\'e de Lausanne (Suisse) \\
2021	& Membre d'un \textbf{Comit\'e de S\'election}, Universit\'e de Franche-Comt\'e \\
		& Membre d'un \textbf{Comit\'e de S\'election}, Universit\'e Autonome de Barcelone (Espagne) \\
2020	& Membre d'un \textbf{Comit\'e de S\'election}, Universit\'e de Paris-Est Cr\'eteil \\
		& Membre d'un \textbf{Comit\'e de S\'election}, Universit\'e de Strasbourg \\
		& Membre d'un \textbf{Comit\'e de S\'election}, Universit\'e de Bourgogne \\
2019	& Membre du jury de \textbf{concours de Directeur de Recherche} de 2\`eme classe "Agronomie, biologie et am\'elioration des plantes, sciences du num\'erique, sciences \'economiques et sociales" de l'INRAE \\
		& Membre d'un \textbf{Comit\'e de S\'election}, Universit\'e de Paris-Est Cr\'eteil \\
2017 	& Membre d'un \textbf{Comit\'e de S\'election} \\
2015 	& Pr\'esidente d'un \textbf{Comit\'e de S\'election}, Universit\'e de Franche-Comt\'e \\	
2015	& Membre de deux \textbf{Comit\'es de S\'election}, Universit\'e de Lorraine et Universit\'e de Lille 3 \\
2014 	& Membre d'un \textbf{Comit\'e de S\'election}, Universit\'e de Strasbourg \\
2013 	& Membre d'un \textbf{Comit\'e de S\'election}, Universit\'e de Bourgogne et  Universit\'e de Franche-Comt\'e \\
2012 	& Membre de deux \textbf{Comit\'es de S\'election}, Universit\'e de Bourgogne \\
2010 	& Membre d'un \textbf{Comit\'e de S\'election}, Universit\'e de Franche-Comt\'e \\
2009 	& Membre de deux \textbf{Comit\'es de S\'election} : Universit\'e de Paris-Sud, Universit\'e d'Orl\'eans \\
2007 - 2008 & Membre suppl\'eant \'elu de la \textbf{Commission de Sp\'ecialistes 5\ieme{} section} de l'U.F.R. de Sciences Juridiques, Economiques, Politiques et de Gestion \\
2007 - 2008 & Membre suppl\'eant \'elu de la \textbf{Commission de Sp\'ecialistes 5\ieme{} section} \`a l'U.F.R. de Sciences Economiques et de Gestion de l'Universit\'e de Bourgogne \\
2005 - 2006 & Membre titulaire nomm\'e de la \textbf{Commission de Sp\'ecialistes 5\ieme{} section} \`a l'U.F.R. de Sciences Economiques et de Gestion de l'Universit\'e de Bourgogne
\end{longtable}

\smalltitre{Autres}

\noindent \begin{longtable}{@{}p{\annee}p{\texte}@{}}
2021 - 		& Membre du \textbf{Conseil Scientifique de Meilleurs Agents} \\
2021 - 		& Membre du \textbf{Comit\'e d'orientation sur la Carte des Loyers - ANIL} \\
& Membre du \textbf{Comit\'e de pilotage du R\'eseau REFINE (Real Estate Finance and Economics Network)} \\
2020 - 		& Membre du \textbf{Conseil Scientifique du R\'eseau REFINE (Real Estate Finance and Economics Network)} \\
2016 - 2020	& Membre nomm\'e du \textbf{Conseil Scientifique du d\'epartement SAE2 de l'INRA devenu d\'epartement Eco-Socio de l'INRAE en 2020} \\
2010 - 2014	& Membre titulaire nomm\'e de la \textbf{Commission r\'egionale d'inscription} et de la \textbf{Chambre r\'egionale de discipline des commissaires aux comptes du ressort de la cour d'appel de Besan\c{c}on} 
\end{longtable}

\newpage


%%% Activit\'es de recherche %%%

\bigtitre{Activit\'es de recherche}

%%% Th\'emes de recherche %%%

\titre{Domaines de recherche actuels}

%\smalltitre{Statistique et \'econom\'etrie spatiales} \medskip \\
%	\ding{219}~\textit{Statistique spatiale exploratoire}: analyse exploratoire des donn\'ees spatiales, cha\'enes de %Markov spatiale, statistique exploratoire spatio-temporelle \medskip \\ 
%	\ding{219}~\textit{Estimation et inf\'erence en \'econom\'etrie spatiale} : mod\'eles en coupe transversale et en panel, %endog\'en\'eit\'e dans les mod\'eles spatiaux, simulations, m\'ethodes de filtrage  \medskip \\
%	\ding{219}~\textit{R\'eseaux spatiaux} \\

%\smalltitre{Economie g\'eographique} \medskip \\
%	\ding{219}~\textit{Prix fonciers et immobiliers} : mod\'eles h\'edoniques, d\'eterminants des prix du foncier et de %l'immobilier, capitalisation des am\'enit\'es environnementales, prix des parcelles de vigne \medskip \\ 
%	\ding{219}~\textit{Recomposition des espaces urbains} : suburbanisation des activit\'es \'economiques, \'etalement %urbain, s\'egr\'egation socio-spatiale \medskip \\ 
%	\ding{219}~\textit{Disparit\'es r\'egionales au sein de l'Union Europ\'eenne}: croissance et politiques r\'egionales, %impact des fonds structurels, convergence des r\'egions europ\'eennes, changement structurel \medskip \\ 
%	\ding{219}~\textit{Autres th\'emes d'\'economie g\'eographique} : d\'eterminants des changements d'usage du sol, %syst\'emes de ville, mesure des \'economies d'agglom\'eration, concurrence fiscale inter-communale, localisation des IDE \\

%\smalltitre{Autres th\'emes de recherche en \'econom\'etrie} \medskip \\
%	\ding{219}~\'evaluation des politiques publiques, mod\'eles hi\'erarchiques, \'econom\'etrie des donn\'ees %exp\'erimentales, m\'eta-analyse, r\'egressions quantiles \medskip \\ 

\smalltitre{Th\`emes} \smallskip

\noindent \textbf{Economie g\'eographique} : mobilit\'es r\'esidentielles, localisation des activit\'es, disparit\'es territoriales, am\'enagement du territoire \smallskip \\ 
\textbf{Economie fonci\`ere et du logement} : usage et artificialisation des sols, d\'eterminants des prix du foncier et de l'immobilier, politiques du logement \smallskip \\ 
\textbf{Economie publique locale} : d\'epenses publiques locales, coop\'eration inter-municipale, p\'er\'equation \medskip \\ 

\smalltitre{M\'ethodes} \smallskip

\noindent \textbf{Micro\'econom\'etrie} : Mod\`{e}les hi\'erarchiques, mod\`{e}les d'\'evaluation d'impact \smallskip \\ 
\textbf{Statistique et \'econom\'etrie sur donn\'ees spatiales} : analyse exploratoire des donn\'ees spatiales, mod\`{e}les \'econom\'etriques spatiaux \smallskip \\ 
\textbf{M\'ethodes exp\'erimentales} : Exp\'eriences en laboratoire, exp\'eriences contr\^ol\'ees, testing \smallskip \\ 

%%% Publications %%%

\titre{Publications}

%%% Revues \'e comit\'e de lecture %%%

\smalltitre{Publications dans des revues \`a comit\'e de lecture} \medskip

\noindent \textit{112 publications dans des revues \`a comit\'e de lecture dont 96 dans des revues class\'ees par la section 05 (\'economie) du CNU (liste de novembre 2025) : 48 en cat\'egorie A, 30 en cat\'egorie B, 18 en cat\'egorie C} \medskip

\noindent \textbf{A para\^itre}  \smallskip \\ 
- Urban Economics for Sustainable Transitions: Grand Challenges in Urban Structure, the Urban Rural Gradient and City Systems, \textit{\textbf{Frontiers in Sustainable Cities Urban Economics}} \\
- From sharing to capitalizing: Assessing the rise of Airbnb in housing prices, \textit{\textbf{Journal of Regional Science}} (avec Y.~Allam et M. Breuill\'e)(\textit{HCERES A}) \\
- Impact of migration status on employment quality: Causal empirical evidence for Tunisian migrants, \textit{\textbf{Applied Economics}} (avec F.~Ben Youssef et M. Kriaa)(\textit{HCERES A}) \\

\noindent \textbf{2026}  \smallskip \\ 
- Econom\'etrie des donn\'ees spatiales : enjeux d'identification et perspectives, \textit{\textbf{R\'egion et D\'eveloppement}}, (avec N. Debarsy), n 62-2025, pp. 245-251 (\textit{HCERES C}) \\
- Is compulsory inter-municipal cooperation an efficiency booster?, \textit{\textbf{Regional Science and Urban Economics}}, vol. 117, pp. 104185 (avec F. Galli, J. Piedra-Pena et M. Breuill\'e)(\textit{HCERES A}) \\
- Can public investment in transport influence densification and land use? Evidence from the Tramway of Dijon (France), \textit{\textbf{Transport Policy}}, vol. 176, pp. 103927 (avec J. Dub\'e, M. Hilal, M.-P. Champagne et D. Legros)(\textit{HCERES B}) \\

\noindent \textbf{2025}  \smallskip \\ 
- Choosing right Bayesian tools: A comparative study of modern Bayesian methods in spatial econometric models, \textit{\textbf{Econometrics}}, vol. 13(4), pp. 49 (avec Y. Ling)(\textit{HCERES C}) \\
- PARIS2019: The impact of rent control on the Parisian rental market, \textit{\textbf{Journal of Housing Economics}}, vol. 70, pp. 102101 (avec M.~Breuill\'e, Y. Morin, M. Regnaud)(\textit{HCERES A}) \\
- Is there an additional price premium for single-family houses exposed to urban parks? Insights from causal spatio-temporal matching in Qu\'ebec City, \textit{\textbf{Journal of Environmental Management}}, vol. 394, pp. 127417 (avec J.~Dub\'e, C. Chapel, M. Hilal, F. Des Rosiers, M.-P. Champagne)(\textit{HCERES B}) \\
- Identification of spatial spillovers: Do's and Dont's, \textit{\textbf{Journal of Economic Surveys}}, vol. 39(5), pp. 2152-2173  (avec N.~Debarsy)(\textit{HCERES A}) \\
- Food security and weather events: a multidimensional analysis in the Sahel for the period 2000-2016, \textit{\textbf{Journal of Asian and African Studies}}, vol. 60(5), pp. 2972-2993 (avec O. Yobom) \\
- From COVID-19 to "new normal"? Revisiting Parisian short-term rental markets and their dynamics, \textit{\textbf{Current Issues in Tourism}}, vol. 28(13), pp. 2173-2188 (avec Y. Ling et M. Breuill\'e)(\textit{HCERES B}) \\
- Tax competition with intermunicipal cooperation, \textit{\textbf{National Tax Journal}}, vol. 78(1), pp. 5-43 (avec D.~Agrawal and M. Breuill\'e)(\textit{HCERES A}) \\
- Review and challenges in the economic valuation of green spaces, \textit{\textbf{Economics and Business Letters}}, vol. 14(1), pp. 116-138 (avec C.~Chapel and M. Hilal) \\
- Discrimination against people with mental, physical or visual disabilities in the French rental market: Field experiment, \textit{\textbf{Housing Studies}}, vol. 40(1), pp. 116-138 (avec A.~Flage) (\textit{HCERES B}) \\

\noindent \textbf{2024}  \smallskip \\ 
- Municipal efficiency spillovers in France, \textit{\textbf{Annals of Regional Science}}, vol. 73, pp.~2143-2172 (avec J. Piedra-Pena and M. Breuill\'e)(\textit{HCERES A}) \\
- Airbnb des villes et des champs \`{a} l'\'epreuve de la Covid-19, \textit{\textbf{Revue d'Economie Industrielle}}, vol. 186, pp. 9-88  (avec Y. Ling, Y. Allam, M. Breuill\'e et C. Grivault)(\textit{HCERES B}) \\
- An integrated causal framework to evalute uplift value with an example on change in public transport supply, \textit{\textbf{Transportation Research Part E}}, vol. 185, pp. 103500  (avec J. Dub\'e, F. Des Rosiers, D. Legros et M.P. Champagne)(\textit{HCERES A}) \\

\noindent \textbf{2023}  \smallskip \\ 
- Assessing French Departments spending efficiency over time, \textit{\textbf{Applied Economics}}, vol.~55, n. 59, pp. 6939-3964 (avec K. Ayouba et M.-L. Duboz)(\textit{HCERES A}) \\
- Bayesian spatial panel models: A flexible Kronecker error component approach, \textit{\textbf{Letters in Spatial and Resource Sciences}}, vol 16(39) (avec Y. Ling) \\
- On the proper computation of the Hausman test statistic in standard linear panel data models: Some clarifications and new results, \textit{\textbf{Econometrics}}, vol. 11(4) (avec M.-A. S\'en\'egas)(\textit{HCERES C}) \\
- Drivers of PES effectiveness: some evidence from a quantitative meta-analysis, \textit{\textbf{Ecological Economics}}, vol. 210, pp. 107856 (avec V. Bellassen, S. Legras, L. Saint-Cyr et L.~V\'edrine)(\textit{HCERES A}) \\
- Spatial panel simultaneous equations models with error components, \textit{\textbf{Empirical Economics}}, vol. 65, pp. 1149-1196 (avec M.C.O. Amba et T. Mbratana)(\textit{HCERES C}) \\

\noindent \textbf{2022}  \smallskip \\ 
- Residential migration and the Covid-19 crisis: Towards an urban exodus in France?,  \textit{\textbf{Economie et Statistique/Economics and Statistics}}, 536-7, pp. 57-73 (avec M. Breuill\'e et A. Verlhiac)(\textit{HCERES B}) \\
- Specification and estimation of a periodic spatial panel autoregressive model, \textit{\textbf{Journal of Spatial Econometrics}}, 3(1) (avec M.C.O. Amba)  \\ 
- Push and pull factors in Tunisian internal migration: the role of human capital, \textit{\textbf{Growth and Change}}, 53(2), pp.~771-799 (avec R. Laajimi) (\textit{HCERES B}) \\ 
- Climate and agriculture: empirical evidence for agroecological zones and countries for the Sahel, \textit{\textbf{Applied Economics}}, 54(8), pp.~918-936 (avec O.~Yobom)(\textit{HCERES A})\\
- The effects of land price in the peri-urban fringe of Mexico City: environmental amenities for informal land parcels purchasers, \textit{\textbf{Urban Studies}}, 59(1), pp.~222-241  (avec E.T.~Martinez-Jimenez, E. P\'erez-Campuzano et A.A. Ibarra) (\textit{HCERES A}) \\ 

\noindent \textbf{2021}  \smallskip \\ 
- Les b\'en\'efices de la densification, \textit{\textbf{Regards Crois\'es sur l'Economie}}, 2021(28), pp.~170-177 (avec M.L.~Breuill\'e et A.~Bretagnolle)\\
- The spatial dimension of the French private rental markets. Evidence from microgeographic data in 2015,
\textit{\textbf{Environment and Planning B}}, 48(8), pp.~2497-2513 (avec K.~Ayouba, M.~Breuill\'e et C.~Grivault)\\
- Parent isol\'e recherche appartement : un test et une \'evaluation, \textit{\textbf{Population}}, 76(2021/1), pp.~77-114 (avec L. Challe, Y. L'Horty, L. du Parquet et P. Petit) (\textit{HCERES B})\\	
- Aux origines des disparit\'es de d\'epenses des d\'epartements fran\c cais : une analyse empirique (2006-2016), \textit{\textbf{Revue d'Economie Politique}}, 2021/2(132), pp.~223-247 (avec M.L.~Duboz et M. Houser) (\textit{HCERES A})\\
- Environmental expenditure interactions among OECD countries, 1995-2017 , \textit{\textbf{Economic Modelling}}, 94, pp.~244-255 (avec Y. NDiaye) (\textit{HCERES A})\\
- Social preferences across different populations: meta-analysis on the ultimatum and dictator games, \textit{\textbf{Journal of Behavioral and Experimental Economics}}, 9, pp.~101613 (avec F. Cochard, N. Georgantzis et J.-C.~Tisserand) (\textit{HCERES B}) \\ 

\noindent \textbf{2020}  \smallskip \\ 
- What geographical concentration of industries in Tunisian Sahel? Empirical evidence using continuous measures, \textit{\textbf{Tijdschrift voor economische en sociale geografie}}, 111(5), pp.~738-757 (avec R. Laajimi et S.~Benammou) \\	
- Faciliter la mobilit\'e quotidienne des jeunes \'eloign\'es de l'emploi : une \'evaluation exp\'erimentale, \textit{\textbf{Revue d'Economie Politique}}, 130(4), pp.~519-544 (avec D. Anne et Y. L'Horty) (\textit{HCERES A})\\
- Regional gatekeepers, inventor networks and inventive performance, \textit{\textbf{Research Policy}}, 49(5), pp.~103981 (avec A. Plunket) (\textit{HCERES A})\\	
- Discrimination in access to housing: a test on urban areas in metropolitan France, \textit{\textbf{Economie et Statistique/Economics and Statistics}}, 513, pp.~27-45 (avec Y. L'Horty, L.~du Parquet et P. Petit) (\textit{HCERES B})\\
- Testing for spatial group-wise heteroskedasticity in spatial autocorrelation regression models: Lagrange multiplier scan tests, \textit{\textbf{Annals of Regional Science}}, 64(2), pp.~287-312 (avec F. L\'opez et C. Chasco) (\textit{HCERES A})\\
- Does Airbnb disrupt the private rental market? An empirical analysis for French cities, \textit{\textbf{International Regional Science Review}}, 43(1-2), pp.~76-104 (avec K.~Ayouba, M.~Breuill\'e et C.~Grivault) (\textit{HCERES~A}) \\		
- Beyond GDP: an analysis of the socio-economic diversity of European regions, \textit{\textbf{Applied Economics}}, 52(9), pp.~1010-1029 (avec K. Ayouba et A. Vallone) (\textit{HCERES A})\\

\noindent \textbf{2019}  \smallskip \\ 
- Gender differences in legal disputes: the case of French labor courts, \textit{\textbf{Revue Economique}}, 70(6), pp.~1201-1212 (avec N. Georgantzis, E.~Peterle et J.-C. Tisserand) (\textit{HCERES A})\\	
- Which attracts more FDI in services: national or regional performance? An empirical analysis of the EU for 1997-2012 , 70(3), pp.~345-373 (avec M.-L. Duboz and N. Kroichvili), \textit{\textbf{Annals of Regional Science}} (\textit{HCERES A})\\ 
- H\'et\'erog\'en\'eit\'e spatiale des prix h\'edoniques des appartements du march\'e locatif priv\'e en France, \textit{\textbf{Revue Fran\c{c}aise d'Economie}} (avec K.~Ayouba, M.~Breuill\'e, C.~Grivault et I.~Nappi-Choulet), XXXIV(3), pp.~203-248 (\textit{HCERES~B}) \\
- Impact de la densification sur les co\^uts des infrastructures et services publics, \textit{\textbf{Revue Economique}}, 63(3), pp.~601-638 (avec M.~Breuill\'e, C.~Grivault et R.~Le Goix) (\textit{HCERES~A}) \\
		
\noindent \textbf{2018}  \smallskip \\ 
- Does issuing building permits reduce the cost of land? An estimation based on the demand for building land, \textit{\textbf{Economie et Statistique/Economics and Statistics}}, 500-501-502, pp.~45-67 (avec J.-S.~Ay, M.~Hilal et J.~Cavailh\`{e}s)(\textit{HCERES~B}) \\
- What role for human capital in the growth process: new evidence from endogenous latent factor panel quantile regression, \textit{\textbf{Scottish Journal of Political Economy}}, 65(5), pp.~501-527 (avec P.~Kostov) (\textit{HCERES~B}) \\
- A multi-dimensional spatial lag panel data model with spatial moving average nested random effects errors, \textit{\textbf{Empirical Economics}}, 55, pp.~113-146 (avec B.~Fingleton et A.~Pirotte) (\textit{HCERES~C}) \\
- Panel data models with spatially dependent nested random effects, \textit{\textbf{Journal of Regional Science}}, 58(1), pp.~63-80 (avec B.~Fingleton et A.~Pirotte) (\textit{HCERES~A}) \\
- Nonlinear impact estimation in spatial autoregressive models, \textit{\textbf{Economics Letters}}, 163, pp.~59-64 (avec K.~Ayouba et J.-S.~Ay) (\textit{HCERES~B}) \\
- A Scan test for spatial groupwise heteroscedasticity in cross-sectional models with an application on house prices in Madrid, \textit{\textbf{Regional Science and Urban Economics}}, 68 pp.~226-238 (avec C.~Chasco et F.~L\'opez) (\textit{HCERES~A}) \\
- Spatial variation in energy attitudes and perceptions: evidence from Europe, \textit{\textbf{Renewable $\&$ Sustainable Energy Reviews}}, 81 pp.~2160-2180 (avec N.~Balta-\"Ozkan) \\
		
\noindent \textbf{2017}  \smallskip \\ 
- Spatial fiscal interactions among French municipalities within inter-municipal groups, \textit{\textbf{Applied Economics}}, 46, pp.~4617-4637 (avec M.-L.~Breuill\'e) (\textit{HCERES~A}) \\
- Assessment of Latin American sustainability, \textit{\textbf{Renewable $\&$ Sustainable Energy Reviews}}, 78, pp.~878-885 (avec O.~Toumi et J. Ben Rejeb) \\
- Aggregated versus individual land-use models: Modeling spatial autocorrelation to increase predictive accuracy, \textit{\textbf{Environmental Modeling and Assessment}}, 22, pp.~129-145 (avec J.-S.~Ay et R.~Chakir) (\textit{HCERES~A}) \\
- Does enhanced mobility of young people improve employment and housing outcomes? Evidence from a large controlled experiment in France, \textit{\textbf{Journal of Urban Economics}}, 97, pp.~1-14 (avec Y.~L'Horty et P.~Petit) (\textit{HCERES~A}) \\
		
\noindent \textbf{2016}  \smallskip \\ 
 - Do foreign investors' location determinants in service functions differ according to sectors? An empirical analysis of EU for 1997-2011, \textit{\textbf{International Regional Science Review}}, 39(4), pp.~417-456 (avec M.-L.~Duboz et N.~Kroichvili) (\textit{HCERES~A}) \\

 
\noindent \textbf{2015}  \smallskip \\ 
 - Regulation of pollution in the laboratory: random inspections, ambient inspections and commitment problems, \textit{\textbf{Bulletin of Economic Research}}, 67(S1), pp.~S40-S73 (avec F.~Cochard et L. Franckx) (\textit{HCERES~B}) \\
- Secret versus public reserve price in an "outcry" English procurement auction: experimental results, \textit{\textbf{International Journal of Production Economics}}, 169, pp.~285-298 (avec K.~Brisset et F.~Cochard)(\textit{HCERES~A}) \\
- Convergence: a story of quantiles and spillovers, \textit{\textbf{Kyklos}}, 68(4), pp.~552-576 (avec P.~Kostov) (\textit{HCERES~A}) \\
- Heterogeneity in perceptions of noise and air pollution: a spatial quantile approach on the agglomeration of Madrid, \textit{\textbf{Spatial Economic Analysis}}, 10(3), pp.~317-343 (avec C.~Chasco) (\textit{HCERES~B}) \\
- Exploring Scan methods to test spatial structure with an application to housing prices in Madrid, \textit{\textbf{Papers in Regional Science}}, 94(2), pp.~316-346 (avec F.~L\'opez et C. Chasco) (\textit{HCERES~A}) \\

\noindent \textbf{2014}  \smallskip \\ 
- Autocorr\'elation spatiale des erreurs et erreurs de mesure : quelles interactions ?, \textit{\textbf{R\'egion et D\'eveloppement}}, 40-2014, pp.~37-52 (avec J.~Mutl) (\textit{HCERES~C}) \\
- L'apport de la politique de coh\'esion aux villes europ\'eennes, \textit{\textbf{Revue Fran\c caise de Finances Publiques}}, 125, pp.~57-68 (avec M.-L.~Duboz) \\

\noindent \textbf{2013}  \smallskip \\ 
- Using synthetic variables in instrumental variable estimation of spatial series models, \textit{\textbf{Environment and Planning A}}, 45(9), pp.~2227-2242 (avec A.~P\'aez) (\textit{HCERES~A}) \\
- Politique urbaine et politique de coh\'esion de l'UE: un lien \`a explorer, \textit{\textbf{Economie Appliqu\'ee}}, LXVI(2), pp.~133-160 (avec M.-L.~Duboz) (\textit{HCERES~C}) \\
- Predicting land use allocation in France: a spatial panel analysis, \textit{\textbf{Ecological Economics}}, 92, pp.~114-125 (avec R.~Chakir) (\textit{HCERES~A}) \\
- The impact of objective and subjective measures of air quality and noise on house prices: a multilevel approach for downtown Madrid, \textit{\textbf{Economic Geography}}, 89(2), pp.~127-148 (avec C.~Chasco) (\textit{HCERES~A}) \\
- The leading role of manufacturing in China's regional economic growth: a spatial econometric approach of Kaldor's laws, \textit{\textbf{International Regional Science Review}}, 36(2), pp.~139-166 (avec D.~Guo et S.~Dall'erba) (\textit{HCERES~A}) \\

\noindent \textbf{2012}  \smallskip \\ 
- Hierarchy and spatial autocorrelation effects in hedonic models, \textit{\textbf{Economics Bulletin}}, 32(2), pp.~1474-1480 (avec C.~Chasco) (\textit{HCERES~B}) \\
 - Autocorr\'elation spatiale et endog\'en\'eit\'e: quelle utilit\'e pour le mod\`ele de Durbin?, \textit{\textbf{Revue d'Economie R\'egionale et Urbaine}}, 2012(1), pp.~3-18 (avec B.~Fingleton) (\textit{HCERES~B}) \\
 - Does manure management regulation work against agglomeration economies? Theory and evidence from the French hog sector, \textit{\textbf{American Journal of Agricultural Economics}}, 94(1), pp.~116-132 (avec C.~Gaign\'e, S.~Larue et B.~Schmitt) (\textit{HCERES~A}) \\
 - Measurement errors in a spatial context, \textit{\textbf{Regional Science and Urban Economics}}, 42(1-2), pp.~114-125 (avec B. Fingleton) (\textit{HCERES~A}) \\
 
\noindent \textbf{2011}  \smallskip \\ 
- Perspectives d'adh\'esion \`a l'UE et \'evolution des \'echanges agricoles des PECO, \textit{\textbf{Economie Rurale}}, 325-326, pp.~54-68 (avec M.-L.~Duboz) (\textit{HCERES~C}) \\
 - The local vs. global dilemna of the effects of structural funds, \textit{\textbf{Growth and Change}}, 42(4), pp.~466-490 (avec S. Dall'erba et R. Guillain) (\textit{HCERES~B}) \\
 - Are EU-15 and CEEC agricultural exports in competition? Evidence for 1995-2005, \textit{\textbf{Economics Bulletin}} (avec M.-L.~Duboz), 31(1), pp.~134-146 (\textit{HCERES~B}) \\
 - The evolution of regional productivity disparities in the European Union from 1975 to 2002: a combination of shift-share and spatial econometrics, \textit{\textbf{Regional Studies}}, 45(1), pp.~123-139 (avec Y.~Kamarianakis) (\textit{HCERES~A}) \\

\noindent \textbf{2010}  \smallskip \\ 
- Agglomeration and dispersion of economic activities in and around Paris: An exploratory spatial data analysis, \textit{\textbf{Environment and Planning B}}, 37(6), pp.~961-981 (avec R.~Guillain) \\

\noindent \textbf{2009}  \smallskip \\ 
- Un regard nouveau sur les politiques de d\'eveloppement r\'egional en Europe , \textit{\textbf{Canadian Journal of Regional Science}}, 32(3), pp.~465-480 (avec S.~Dall'erba et R.~Guillain) (\textit{HCERES~C}) \\
 - Impact of structural funds on regional growth: How to reconsider a 9-year-old black box, \textit{\textbf{R\'egion et D\'eveloppement}}, 2009-30, pp.~77-100 (avec S.~Dall'erba et R.~Guillain) (\textit{HCERES~C}) \\
 - The state-federal dichotomy in the effects of minimum wages on teenage employment in the United States, \textit{\textbf{Economics Letters}}, 105(3), pp.~267-269 (avec S.~Bazen) (\textit{HCERES~B}) \\
 - Les sch\'emas de concentration sectorielle au sein de l'Union Europ\'eenne : l'Est miroir de l'Ouest ?, \textit{\textbf{Economie et Statistique}}, 423, pp.~59-76 (avec M.-L.~Duboz et R.~Guillain) (\textit{HCERES~B}) \\
		
\noindent \textbf{2008}  \smallskip \\ 
- Spatial and sectoral productivity convergence between European regions, 1975-2000, \textit{\textbf{Papers in Regional Science}}, 87(4), pp.~505-525 (avec S.~Dall'erba) (\textit{HCERES~A}) \\
 - Estimating spatial models with endogenous variables, a spatial lag and spatially dependent disturbances: finite sample properties, \textit{\textbf{Papers in Regional Science}}, 87(3), pp.~319-339 (avec B.~Fingleton) (\textit{HCERES~A}) \\
- Does evidence on regional economic convergence depend on the estimation strategy? Outcomes from analysis of a set of NUTS-2 regions, \textit{\textbf{Spatial Economic Analysis}}, 3(2), pp.~209-224 (avec G.~Arbia et G.~Piras) (\textit{HCERES~B}) \\
- Regional convergence and the impact of European structural funds over 1989-1999: a spatial econometric analysis, \textit{\textbf{Papers in Regional Science}}, 87(2), pp.~219-244 (avec S.~Dall'erba) (\textit{HCERES~A}) \\
- Spatial analysis of urban growth in Spain, 1900-2001, \textit{\textbf{Empirical Economics}}, 34(1), pp.~59-80 (avec C.~Chasco) (\textit{HCERES~C}) \\
- Le centre d'affaire historique de Paris : quel pouvoir structurant sur l'espace \'economique en Ile-de-France ?, \textit{\textbf{Revue d'Economie Industrielle}}, 123, pp.~65-86 (avec R.~Guillain) (\textit{HCERES~B}) \\
 - Identifier la localisation des activit\'es \'economiques : une approche par les outils de l'analyse exploratoire des donn\'ees spatiales, \textit{\textbf{Economie Appliqu\'ee}}, LXI(3), pp.~5-34 (avec R.~Guillain) (\textit{HCERES~C}) \\
- Fonds structurels, effets de d\'ebordement g\'eographique et croissance r\'egionale en Europe, \textit{\textbf{Revue de l'OFCE}}, 104, pp.~241-270 (avec S.~Dall'erba et R.~Guillain) (\textit{HCERES~C}) \\

\noindent \textbf{2007}  \smallskip \\ 
- Finite sample properties of estimators of spatial models with autoregressive, or moving average, disturbances and system feedback, \textit{\textbf{Annales d'Economie et de Statistique}}, 87-88, pp.~39-62 (avec B.~Fingleton) (\textit{HCERES~A}) \\	
 - Local versus global convergence in Europe: a bayesian spatial econometric approach, \textit{\textbf{Review of Regional Studies}}, 37(1), pp.~82-108 (avec C.~Ertur et J.~LeSage) \\
 - The impact of EU regional support on growth and employment, \textit{\textbf{Czech Journal of Economics and Finance - Finance a \'uv\v{e}r}}, 57(7-8), pp.~325-340 (avec S.~Dall'erba) (\textit{HCERES~C}) \\
 - La valeur \'economique des paysages des villes p\'eriurbanis\'ees, \textit{\textbf{Economie Publique}}, 20(1), pp.~3-27 (avec T.~Brossard, J.~Cavailh\`es, G.~G\'eniaux, M.~Hilal, H.~Jayet, D.~Joly, C.~Napoleone, N.~Ovtracht, \hbox{P.-Y.~P\'eguy},\hbox{ F.-P.~Tourneux} et P.~Wavresky) (\textit{HCERES~C}) \\
		
\noindent \textbf{2006}  \smallskip \\ 
- Changes in spatial and sectoral patterns of employment in Ile-de-France, 1978-1997, \textit{\textbf{Urban Studies}}, 43(11), pp.~2075-2098 (avec R.~Guillain et C.~Boiteux-Orain) (\textit{HCERES~A}) \\
 - Interpolation of air quality measures in hedonic house price models: spatial aspects, \textit{\textbf{Spatial Economic Analysis}}, 1(1), pp.~31-52 (avec L.~Anselin) (\textit{HCERES~B}) \\
 - Evaluating the temporal and spatial heterogeneity of the European convergence process: 1980-1999, \textit{\textbf{Journal of Regional Science}}, 46(2), pp.~269-288 (avec S.~Dall'erba) (\textit{HCERES~A}) \\
 - The European regional convergence process, 1980-1995: Do spatial dependence and spatial heterogeneity matter?, \textit{\textbf{International Regional Science Review}}, 29(1), pp.~2-34 (avec C.~Ertur et C.~Baumont) (\textit{HCERES~A}) \\
 - Clubs de convergence et effets de d\'ebordement g\'eographiques : une analyse spatiale sur donn\'ees r\'egionales europ\'eennes, 1980-1995, \textit{\textbf{Economie et Pr\'evision}}, 173(2), pp.~111-134 (avec C.~Baumont et C.~Ertur) (\textit{HCERES~B}) \\

\noindent \textbf{2005}  \smallskip \\ 
- Regional productivity differentials in three new member countries. What can we learn from the 1986 enlargement to the South?, \textit{\textbf{Review of Regional Studies}}, 35(1), pp.~97-116 (avec S. Dall'erba, Y.~Kamarianakis et M.~Plotnikova) \\
 - On the property of diffusion in the spatial error model, \textit{\textbf{Applied Economics Letters}}, 12(9), pp.~533-536 (avec C.~Baumont, S.~Dall'erba et C.~Ertur) (\textit{HCERES~C}) \\
 - Dynamique du processus de convergence r\'egionale en Europe, \textit{\textbf{R\'egion et D\'eveloppement}}, 2005-21, pp.~119-139 (avec S.~Dall'erba) (\textit{HCERES~C}) \\
		
\noindent \textbf{2004}  \smallskip \\ 
- Spatial analysis of employment and population: the case of the agglomeration of Dijon, 1999, \textit{\textbf{Geographical Analysis}}, 36(2), pp.~146-176 (avec C.~Baumont et C.~Ertur) \\ 
 - Space-time analysis of GDP disparities among European regions: A Markov chains approach, \textit{\textbf{International Regional Science Review}}, 27(2), pp.~138-163 (\textit{HCERES~A}) \\
 - La dynamique des disparit\'es r\'egionales dans l'Union Europ\'eenne, \textit{\textbf{Revue d'Economie R\'egionale et Urbaine}}, 4, pp.~491-512 (\textit{HCERES~B}) \\
 - H\'et\'erog\'en\'eit\'e spatiale, principes et m\'ethodes, \textit{\textbf{Economie et Pr\'evision}}, 162(1), pp.~151-172 (\textit{AERES~B}) \\

\noindent \textbf{2003}  \smallskip \\ 
- Exploratory spatial data analysis of the distribution of regional per capita GDP in Europe, 1980-1995, \textit{\textbf{Papers in Regional Science}}, 82(2), pp.~175-201 (avec C.~Ertur) (\textit{HCERES~A}) \\
 - Les diff\'erentiels de productivit\'e r\'egionale dans les pays en transition par rapport \`a la moyenne europ\'eenne : le cas de la Pologne, de la Hongrie et de la R\'epublique Tch\'eque, \textit{\textbf{R\'egion et D\'eveloppement}}, 2003-18, pp.~111-129 (avec S.~Dall'erba, Y.~Kamarianakis et M.~Plotnikova) (\textit{HCERES~C}) \\

\noindent \textbf{2002}  \smallskip \\ 
- Econom\'etrie spatiale : l'autocorr\'elation spatiale dans les mod\'eles de r\'egression lin\'eaire, \textit{\textbf{Economie et Pr\'evision}}, 155(4), pp.~139-158 (\textit{HCERES~B}) \\
 - Convergence des r\'egions europ\'eennes, une approche par l'\'econom\'etrie spatiale, \textit{\textbf{Revue d'Economie R\'egionale et Urbaine}}, 2, pp.~203-216 (avec C.~Baumont et C.~Ertur) (\textit{HCERES~B}) \\

%%% Num\'eros sp\'eciaux %%%

\smalltitre{Editoriaux - Coordination de num\'eros sp\'eciaux}\medskip

\noindent - Introduction of Julie Le Gallo, Editor-in-chief as of January 2025, \textit{\textbf{Journal of Geographical Systems}}, 2025, 27(5-6), pp.~5-6 \\ 
- Recent developments in spatial statistics and spatial econometrics, \textit{\textbf{Annals of Regional Science}}, 2020, 64(2), pp.~239-241 (avec R.~Guillain) \\ 
 - La structuration de l'espace \'economique. Quels apports d'une approche multidisciplinaire ?, \textit{\textbf{Revue Economique}}, 2019, 70(3), pp.~301-304 (avec F.~Puech) \\ 
 - Analyse \'economique de changement d'usage du sol. Introduction, \textit{\textbf{Revue Economique}}, 2017, 68, pp.~405-408 (avec R.~Guillain) \\ 
- Spatial econometrics principles and challenges in Jean Paelinck's research, \textit{\textbf{Spatial Economic Analysis}}, 2015, 10(3), pp.~263-269 (avec C.~Chasco) \\ 
- Statistique, \'econom\'etrie et espace : progr\`{e}s et d\'efis en science r\'egionale, \textit{\textbf{Revue d'Economie R\'egionale et Urbaine}}, 2015, 2015 - 1/2, pp. 9-23 (avec I.~Thomas) \\ 
 - Introduction to special issue: Contributions to spatial econometrics, \textit{\textbf{International Regional Science Review}}, 2014, 37(3), pp.~247-250 (avec R.J.G.M.~Florax et Y.~Bao) \\ 
- Introduction to special issue on revisiting convergence, \textit{\textbf{Economics and Business Letters}}, 2013, 2(4), pp.~140-142 (avec F.J.~Delgado) \\
- Introduction au num\'ero sp\'ecial : Croissance et gouvernance r\'egionale en Europe, \textit{\textbf{R\'egion et D\'eveloppement}}, 2009, 2009-30, pp.~5-10 (avec R.~Guillain) \\
- Introduction au num\'ero sp\'ecial : Economie industrielle et \'econom\'etrie spatiale, \textit{\textbf{Revue d'Economie Industrielle}}, 2008, 123, pp.~15-18 (avec N.~Massard et A.~Plunket) \\
- Introduction au num\'ero sp\'ecial : Croissance, convergence et interactions spatiales, application des outils r\'ecents de l'analyse spatiale quantitative, \textit{\textbf{R\'egion et D\'eveloppement}}, 2005, 2005-21, pp.~5-11 (avec S.~Dall'erba) \\

%%% Ouvrages %%%

\smalltitre{Livre} \medskip

\noindent - \textit{\textbf{Evaluation d'impact des politiques publiques d'am\'enagement. Une approche pratique bas\'ee sur le march\'e immobilier}}, 2024, De Boeck (avec J.~Dub\'e, N.~Devaux et D.~Legros) \\

\smalltitre{Editeur d'un ouvrage collectif} \medskip

\noindent - \textit{\textbf{Progress in Spatial Analysis: Theory and Computation, and Thematic Applications}}, 2009, Advances in Spatial Sciences, Springer-Verlag, Berlin (avec A.~P\'aez, S.~Dall'erba et R.~Buliung) \\

%%% Chapitres ouvrages %%%

\smalltitre{Publications dans des chapitres d'ouvrage collectifs}\medskip

\noindent - Local public cooperation, dans: Nijkamp P., Kourtit K., Haynes K., Elburz Z. (Eds), \textit{\textbf{Thematic Encyclopedia of Regional Science}}, Edward Elgar, 2025, pp.~545-546 (avec M.-L. Breuill\'e). \\
- Fiscal equalization, dans: Nijkamp P., Kourtit K., Haynes K., Elburz Z. (Eds), \textit{\textbf{Thematic Encyclopedia of Regional Science}},  Edward Elgar, 2025, pp.~501-502 (avec M.-L. Breuill\'e) \\
- The spatial structure of housing process in Madrid: Evidence from spatio-temporal scan statistics, dans : Glaz~J., Koutras~M.V. (Eds.), \textit{\textbf{Handbook of Scan Statistics}}, Springer, 2024, pp.~683-701 (avec C.~Chasco et F.~Lopez) \\
- Carte et mod\'eles statistique pour explorer l'h\'et\'erog\'en\'eit\'e spatiale, dans : Cunty~C., Mathian~H. (Eds.), \textit{\textbf{Traitements et cartographie pour l'information g\'eographique}}, 2023, Springer, pp.~147-188 (avec M.~Hilal) \\
- Does cohesion policy affect regional development and (intra-)regional disparities?, dans : Rauhut~D., Sielker~F., Humer~A. (Eds.), \textit{\textbf{EU Cohesion Policy and Spatial Governance}}, 2021, Edward Elgar, pp.~156-170  (avec L.~V\'edrine) \\
- Spatial autocorrelation in econometric land use models: an overview, dans : Daouia~A., Ruiz-Gazen~A. (Eds.), \textit{\textbf{Advances in Contemporary Statistics and Econometrics}}, 2021, Springer, pp. 339-362 (avec R.~Chakir) \\
- Endogeneity in spatial models, dans : Fischer~M., Nijkamp~P. (Eds.), \textit{\textbf{Handbook of Regional Science, 2nd edition}}, 2021, Springer-Verlag, pp. 2237-2265  (avec B.~Fingleton) \\
- Cross-section spatial regression models, dans : Fischer~M., Nijkamp~P. (Eds.), \textit{\textbf{Handbook of Regional Science, 2nd edition}}, 2021, Springer-Verlag, pp. 2117-2139  \\
- Regional growth and convergence empirics, dans : Fischer~M., Nijkamp~P. (Eds.), \textit{\textbf{Handbook of Regional Science, 2nd edition}}, 2021, Springer-Verlag, pp. 679-706 (avec B.~Fingleton) \\
- Heterogeneous reaction versus interaction in spatial econometric regional growth and convergence models, dans : Capello~R., Nijkamp P. (Eds.) \textit{\textbf{Handbook of Regional Growth and Development Theories, 2nd edition}}, 2019, Edwar Elgar (avec C. Ertur) \\
- Les effets de la densit\'e sur les co\^uts des infrastructures et services publics, dans : Prager J.-C. (Ed.), \textit{\textbf{Les effets \'economiques et urbains du Grand Paris Express}}, 2019, Economica, Paris (avec M. Bart\'el\'emy, A. Br\`es, A.~Bretagnolle, M.-L. Breuille, C. Grivault, R. Le Goix et V. Volpati) \\
- Econom\'etrie spatiale sur donn\'ees de panel, dans : de Bellefon~M.-P., Loonis V. (Eds.), \textit{\textbf{Manuel d'Analyse Spatiale}}, 2018, INSEE M\'ethodes n. 131, pp. 185-212 (avec S.~Bouayad Agha et L.~V\'edrine) \\
 - Models for spatial panels, dans : Matyas~L., (Eds.), \textit{\textbf{The Econometrics of Multi-Dimensional Panels - Theory and Application}}, 2017, Springer-Verlag, pp. 263-289 (avec A.~Pirotte) \\
- Cross-section spatial regression models, dans : Fischer~M., Nijkamp~P., (Eds.), \textit{\textbf{Handbook of Regional Science}}, 2014, Springer-Verlag, pp.~291-315 \\
- Regional growth and convergence empirics, dans : Fischer~M., Nijkamp~P., (Eds.), \textit{\textbf{Handbook of Regional Science}}, 2014, Springer-Verlag, pp.~291-315 (avec B.~Fingleton) \\
- Dealing with data at various spatial scales and supports: an application to traffic noise and air pollution effects on housing prices with multilevel models, dans : Fern\'andez-Vazquez~E., Rubierra-Moroll\'on~F. (Eds.), \textit{\textbf{Defining the Spatial Scale in Modern Regional Analysis}}, 2012, Springer-Verlag, Berlin, pp.~281-310 (avec C.~Chasco) \\
- Endogeneity in a spatial context: properties of estimators, dans : P\'aez~A., Le Gallo~J., Dall'erba~S. et Buliung~R. (Eds.), \textit{\textbf{Progress in Spatial Analysis: Theory and Computation, and Thematic Applications}}, 2009, Springer-Verlag, Berlin, pp.~59-73 (avec B.~Fingleton) \\
- Employment density in Ile-de-France: evidence from local regressions, dans : P\'aez~A., Le Gallo~J., Dall'erba~S. et Buliung~R. (Eds.), \textit{\textbf{Progress in Spatial Analysis: Theory and Computation, and Thematic Applications}}, 2009, Springer-Verlag, Berlin, pp.~233-251 (avec R.~Guillain) \\
- Spatial analysis of economic convergence, dans : Mills~T.C. et Patterson~K. (Eds.), \textit{\textbf{The Palgrave Handbook of Econometrics Volume II: Applied Econometrics}}, 2009, Palgrave-MacMillan, pp.~1251-1292 (avec S.J.~Rey) \\
- Regional growth and convergence: heterogeneous reaction versus interaction in spatial econometric approaches, dans : Capello~R. et Nijkamp~P. (Eds.), \textit{\textbf{Handbook of Regional Growth and Development Theories}}, 2009, Edward Elgar, pp.~374-388 (avec C.~Ertur) \\
- La dimension spatiale de la s\'egr\'egation, dans : Gaschet~F. et Lacour~C. (Eds.), \textit{\textbf{M\'etropolisation et s\'egr\'egation}}, 2009, Presses Universitaires de Bordeaux, pp.~45-65 (avec F. Gaschet) \\
- Spatial analysis of urban growth in Spain, 1900-2001, dans : Arbia~G. et Baltagi~B.H. (Eds.), \textit{\textbf{Spatial Econometrics: Methods and Applications}}, 2008, Physica-Verlag, Studies in Empirical Economics, Springer, Berlin, pp. 59-80 (avec C. Chasco) \\
- Spatial econometrics and panel data models, dans : Matyas~L. et Sevestre~P. (Eds.), \textit{\textbf{The Econometrics of Panel Data}}, 2008, Kluwer Academic Publishers, 3\ieme{} edition, pp.~625-660 (avec L.~Anselin et H.~Jayet) \\
- La localisation des activit\'es \'economiques en Ile-de-France et dans les d\'epartements voisins : identification des clusters \`a l'aide d'une analyse exploratoire des donn\'ees spatiales, dans : Larceneux~A. et Boiteux-Orain~C. (Eds.), \textit{\textbf{Paris et ses franges: \'etalement urbain et polycentrisme}}, 2006, Editions Universitaires de Dijon, pp.~203-230 (avec R.~Guillain) \\
- Convergence spatiale et sectorielle de la productivit\'e du travail en Europe, dans : Capron~H. (Ed.), \textit{\textbf{Convergence et dynamique d'innovation au sein de l'Espace Europ\'een}}, 2006, Editions De Boeck, pp.~107-133 (avec S.~Dall'erba) \\
- An exploratory spatial data analysis of European regional disparities, 1980-1995, dans : Fingleton~B. (Ed.), \textit{\textbf{European Regional Growth}}, 2003, Springer-Verlag, Advances in Spatial Science, pp.~55-98 (avec C.~Ertur) \\
- A spatial econometric analysis of convergence across European regions, 1980-1995, dans : Fingleton~B. (Ed.), \textit{\textbf{European Regional Growth}}, 2003, Springer-Verlag, Advances in Spatial Science, pp.~99-130 (avec C.~Ertur et C.~Baumont) \\
- Spatial regimes in the European growth process, 1980-1995, dans : Fingleton~B. (Ed.), \textit{\textbf{European Regional Growth}}, 2003, Springer-Verlag, Advances in Spatial Science, pp.~131-158 (avec C.~Baumont et C.~Ertur) \\
- Les nouvelles centralit\'es urbaines, dans : Baumont~C., Combes~P.-P., Derycke~P.-H. et Jayet~H. (Eds), \textit{\textbf{Economie g\'eographique : les th\'eories \`a l'\'epreuve des faits}}, 2000, Economica, Biblioth\'eque de Science R\'egionale, pp.~211-239 (avec C.~Baumont) \\

%%% Comptes rendus %%%

\smalltitre{Comptes rendus} \medskip

\noindent - New Developments in Spatial Economics and Economic Geography, McCann~P. (Eds), \textit{\textbf{Journal of Regional Science}}, 2016, 56(1), pp.~184-186. \\
- Regional Economic Policy in Europe. New Challenges for Theory, Empirics and Normative Interventions, Stierle-Von Sch\'etz~H., Stierle~M.H., Jennings Jr~F.B., Kuah A.T.H. (Eds.), 2009, \textit{\textbf{R\'egion et D\'eveloppement}}, 2009-30.\\
- Land Use Models, Priemus~H., Button~K., Nijkamp~P. (Eds.), 2008, \textit{\textbf{R\'egion et D\'eveloppement}}, 2008-27. \\
- Planning Models, Reggiani~A., Button~K., Nijkamp~P. (Eds.), 2007, \textit{\textbf{R\'egion et D\'eveloppement}}, 2007-25. \\
- Spatial Dynamics, Networks and Modelling, Reggiani~A., Nijkamp~P. (Eds.), 2007, \textit{\textbf{R\'egion et D\'eveloppement}}, vol.~2007-25. \\

%%% Rapports de recherche %%%

\smalltitre{Rapports de recherche} \medskip

\noindent 
- Vers un nouveau zonage caract\'erisant la tension immobili\`ere, rapport final, \textit{\textbf{Rapport de recherche pour le Minist\`{e}re de la Coh\'esion des Territoires}}, 2025 (avec M.-L.~Breuill\'e, A.~Como, E.~Fiorentino, C.~Grivault et M~Regnaud,) 47 pages \\
- Rapport final du comit\'e scientifique de l'\'evaluation de l'exp\'erimention "Territoires z\'ero ch\^omeur de longue dur\'ee" (TZCLD) - Vers une garantie d'emploi ?, \textit{\textbf{DARES}}, 2025 sous la direction de Y.~Lhorty), 199 pages \\
- D\'emat\'erialisation et acc\`{e}s \`{a} la justice, rapport final, \textit{\textbf{DEMAJUST 2020}}, 2023 (avec M.~Hilal, A.~Diallo, E.~Doidy, E.~Dumont, C.~Tirvaudey et L.~Kondratuk), 377 pages \\
- Evaluation d'impact de l'encadrement des loyers \`{a} PARIS, rapport final, \textit{\textbf{Rapport de recherche pour l'APUR et la Ville de Paris}}, 2023 (avec M.-L.~Breuill\'e, C.~Grivault, Y.~Morin et M.~Regnaud), 90 pages \\
- Impact \'economique des espaces verts \`{a} Qu\'ebec : Investigation sur la base du march\'e r\'esidentiel unifamilial, rapport final, \textit{\textbf{Rapport de recherche pour la ville de Qu\'ebec}}, 2023 (avec J.~Dub\'e, M.~Hilal, C.~Chapel, F.~Des Rosiers et N.~Devaux), 45 pages \\
- Location touristique meubl\'ee : D\'eterminants, impacts et effets de la r\`{e}glementation, rapport final, \textit{\textbf{Rapport de recherche pour le Minist\`{e}re de la Coh\'esion des Territoires}}, 2022 (avec Y.~Allam, M.-L.~Breuill\'e et C.~Grivault), 48 pages \\
- Centralit\'es : comment les identifier et quels r\^oles dans les dynamiques locales et intercommunales ?, \textit{\textbf{Rapport de recherche pour le Commissariat G\'en\'eral \`a l'Egalit\'e des Territoires}}, 2020 (avec H.~Bouscasse, J.~Cavailh\`es, A.~Diallo, M.~Drut, M.~Hilal, S.~Legras, D.~Moret, V.~Piguet), 124 pages \\
- D3.4. Report on the impact of European Cohesion Policy on economic growth and spatial inequalities, \textit{\textbf{IMAJINE H2020 Project}}, 2020 (avec L.~V\'edrine, P. Artelaris et M.~Woods), 105 pages \\
- Les b\'en\'efices de la densification dans la perspective du Grand Paris Express, rapport final, \textit{\textbf{Rapport de recherche pour la Soci\'et\'e du Grand Paris}}, 2020 (avec A.~Bretagnolle, M.-L.~Breuill\'e, C.~Grivault, L.~Guibard, E.~Kakpo, T.~Le Corre, R.~Le Goix, L.~Liliane, L.~Lucas, A.~Pavard), 97 pages \\
- D2.5. Report on spatiotemporal ESDA on GDP, income and educational attainment in European regions , \textit{\textbf{IMAJINE H2020 Project}}, 2019 (avec K.~Ayouba), 75 pages \\
- Convention de recherche relative \`a la d\'etermination d'indicateurs de loyer, rapport final, \textit{\textbf{Rapport de recherche pour le Minist\`{e}re de la Coh\'esion des Territoires}}, 2019 (avec K.~Ayouba, M.-L.~Breuill\'e et C.~Grivault), 79 pages \\
- Convention de recherche relative \`a la d\'etermination d'indicateurs de loyer, rapport interm\'ediaire, \textit{\textbf{Rapport de recherche pour le Minist\`{e}re de la Coh\'esion des Territoires}}, 2018 (avec K.~Ayouba, M.-L.~Breuill\'e et C.~Grivault), 65 pages \\
- Les b\'en\'efices de la densification dans la perspective du Grand Paris Express, livrable n.4, \textit{\textbf{Rapport de recherche pour la Soci\'et\'e du Grand Paris}}, 2018 (avec M.~Barth\'el\'emy, A.~Br\`{e}s, A.~Bretagnolle, M.-L.~Breuill\'e, C.~Grivault, R.~Le Goix, A.~Pavard et V.~Volpati), 188 pages \\
- Parent isol\'e recherche appartement : un test et une \'evaluation (PIRATE), livrable final, \textit{\textbf{Rapport de recherche pour la Mairie de Paris}}, 2018 (avec L.~Challe, Y.~L'Horty, L.~du Parquet et P.~Petit), 19 pages \\
- Deuxi\`{e}me convention de recherche relative aux d\'eterminants de la formation des niveaux de loyer, livrable final, \textit{\textbf{Rapport de recherche pour le Minist\`{e}re du Logement et de l'Habitat Durable}}, 2017 (avec K.~Ayouba, M.-L.~Breuill\'e, M.~Garcia, C.~Grivault et I.~Nappi-Choulet), 158 pages \\
- Deuxi\`{e}me convention de recherche relative aux d\'eterminants de la formation des niveaux de loyer, livrable interm\'ediaire, \textit{\textbf{Rapport de recherche pour le Minist\`{e}re du Logement et de l'Habitat Durable}}, 2017 (avec K.~Ayouba, M.-L.~Breuill\'e, M.~Garcia, C.~Grivault et I.~Nappi-Choulet), 175 pages \\
- Les b\'en\'efices de la densification dans la perspective du Grand Paris Express, livrable n. 3, \textit{\textbf{Rapport de recherche pour la Soci\'et\'e du Grand Paris}}, 2017 (avec M.~Barth\'el\'emy, A.~Br\`{e}s, A.~Bretagnolle, M.-L.~Breuill\'e, C.~Grivault, R.~Le Goix, A.~Pavard et V.~Volpati), 126 pages \\
- Les b\'en\'efices de la densification dans la perspective du Grand Paris Express, livrable n. 2, \textit{\textbf{Rapport de recherche pour la Soci\'et\'e du Grand Paris}}, 2017 (avec M.~Barth\'el\'emy, A.~Br\`{e}s, A.~Bretagnolle, M.-L.~Breuill\'e, C.~Grivault, R.~Le Goix, A.~Pavard et V.~Volpati), 76 pages \\
- Les b\'en\'efices de la densification dans la perspective du Grand Paris Express, livrable n. 1, \textit{\textbf{Rapport de recherche pour la Soci\'et\'e du Grand Paris}}, 2017 (avec M.~Barth\'el\'emy, A.~Br\`{e}s, A.~Bretagnolle, M.-L.~Breuill\'e, C.~Grivault, R.~Le Goix, A.~Pavard et V.~Volpati), 175 pages \\
- Evaluation sociologique et statistique des actions innovantes pour la mobilit\'e et l'accompagnement des jeunes vers l'emploi (ESSAIMAJE), livrable final, \textit{\textbf{Rapport de recherche pour le Fonds d'Exp\'erimentation pour la Jeunesse}}, 2017 (avec D.~Anne et Y.~L'Horty), 60 pages \\
- Convention de recherche relative aux d\'eterminants de la formation des niveaux de loyer, livrable final, \textit{\textbf{Rapport de recherche pour le Minist\`{e}re du Logement et de l'Habitat Durable}}, 2016 (avec K.~Ayouba, M.-L.~Breuill\'e, C.~Grivault et I.~Nappi-Choulet), 446 pages \\
- Convention de recherche relative aux d\'eterminants de la formation des niveaux de loyer, deuxi\'eme livrable, \textit{\textbf{Rapport de recherche pour le Minist\`{e}re du Logement et de l'Habitat Durable}}, 2016 (avec K.~Ayouba, M.-L.~Breuill\'e, C.~Grivault et I.~Nappi-Choulet), 95 pages \\
- Convention de recherche relative aux d\'eterminants de la formation des niveaux de loyer, premier livrable, \textit{\textbf{Rapport de recherche pour le Minist\`{e}re du Logement et de l'Habitat Durable}}, 2015 (avec K.~Ayouba, M.-L.~Breuill\'e, C.~Grivault et I.~Nappi-Choulet), 63 pages \\

%%% actes de colloque %%%

\smalltitre{Actes de colloque} \medskip

\noindent - Is geographical indication acting on rice export price?, \textit{\textbf{Proceedings in System Dynamics and Innovation in Food Network 2018}} (avec G.~Giraud et H.~Boucher) \\

%%% Autres %%%

\smalltitre{Autres} \medskip

\noindent 
- La tension immobili\`{e}re en France, \textit{\textbf{Territoires et transitions. Enjeux du logement. Cahier de l'Observatoire des Territoires}}, 2026 (avec M.~Breuill\'e, A.~Caumo, E.~Fiorentino, C.~Grivault et M.~Regnaud) \\
- Entretien avec... l'\'equipe POPSU Exode urbain, \textit{\textbf{BelvederR, Revue collaborative de l'AUAT}}, 2022, 10 (avec F.~Baillet, M.~Breuill\'e, C.~Grivault et O.~Lision) \\
- Comment \'evaluer l'impact environnemental des usages des sols?, \textit{\textbf{Les Echos, Le Cercle, 2018}} (avec R.~Guillain) \\
- Les r\'egulations environnementales visant l'\'epandage de lisier limitent-elles la concentration g\'eographique des \'elevages porcins ?, \textit{\textbf{INRAE Sciences Sociales}}, 2013, 3/2013 (avec J.~Abildtrup, C.~Gaign\'e, S.~Larue et B.~Schmitt) \\
- Economies d'agglom\'eration et co\^uts de la concentration. Dynamiques de localisation des syst\`emes d'\'elevage intensifs : le cas de la production porcine , \textit{\textbf{Comptes-Rendus de l'Acad\'emie d'Agriculture}}, colloque Elevages intensifs et enviroannement. Les effluents: menace ou richesse ?, 2010, pp.~49-62 (avec S.~Larue, C.~Gaign\'e, J.~Abildtrup, L.~Latruffe et B.~Schmitt) \\
- Les politiques europ\'eennes de coh\'esion territoriale (Compl\'ement K), contribution au rapport du Conseil d'Analyse Economique : \og\textbf{Innovations et comp\'etitivit\'e des r\'egions}\fg~par T.~Madi\'es et J.-C.~Prager, 2008, 77, pp.~343-366 (avec S.~Riou) \\
%%% Contrats, animation et valorisation de la recherche %%%

\newpage

\titre{Contrats, animation et valorisation de la recherche}

\smalltitre{Direction de contrats de recherche}

\noindent \begin{longtable}{@{}p{\annee}p{\texte}@{}}

2025 - 2026 & \og{}\textit{Les effets du Grand Paris Express sur les dynamiques immobili\`eres, r\'esidentielles et touristiques}\fg{}, Projet de recherche financ\'e par la Soci\'et\'e des Grands Projets, co-direction avec M.-L. Breuill\'e (CESAER) \\
2024 - 2025 & \og{}\textit{Indicateurs de tension immobili\`eres en France}\fg{}, Projet de recherche financ\'e par le Minist\`ere de la Coh\'esion des Territoires et des Relations avec les Collectivit\'es Territoriales, co-direction avec M.-L. Breuill\'e (CESAER) \\
2023 - 2026 & \og{}\textit{ArtiFis - Artificialisation des sols en France : d\'eterminants, modes de r\'egulation et cons\'equences sur les march\'es immobiliers}\fg{}, Projet de recherche financ\'e par la R\'egion Bourgogne Franche-Comt\'e, co-direction avec M.-L.~Breuill\'e (CESAER) \\
2023 - 2025 & \og{}\textit{Des mythes m\'ediatiques aux r\'ealit\'es socio-\'economiques : dynamiques spatiales et h\'et\'erog\'en\'eit\'e des espaces ruraux}\fg{}, Projet de recherche financ\'e par le PUCA, co-direction avec M.-L. Breuill\'e (CESAER) \\
2022 - 2023 & \og{}\textit{Evaluation de la performance p\'er\'equatrice des transferts de l'Etat aux communes et EPIC}\fg{}, Projet de recherche dans le cadre des travaux du r\'eseau FInances locales, co-direction avec M.-L. Breuill\'e (CESAER) \\
2022 - 2024 & \og{}\textit{Evaluation d'impact de l'encadrement des loyers \`a Paris}\fg{}, Projet de recherche financ\'e par l'APUR, co-direction avec M.-L. Breuill\'e (CESAER) \\
2021 - 2022 & \og{}\textit{Comment la plateforme d'\'economie collaborative Airbnb modifie-t-elle le march\'e du logement? }\fg{}, Projet de recherche financ\'e par le Minist\`ere de la Coh\'esion des Territoires et des Relations avec les Collectivit\'es Territoriales, co-direction avec M.-L. Breuill\'e (CESAER) \\
& \og{}\textit{Pendant la crise sanitaire, observer les envies de mobilit\'e et leurs concr\'etisations, gr\^ace \`a la digitalisation de l'immobilier}\fg{}, Projet de recherche financ\'e par le PUCA, co-direction avec M.-L. Breuill\'e (CESAER) \\
2018 - 2020 & \og{}\textit{Indicateurs de loyers}\fg{}, Projet de recherche financ\'e par le Minist\`ere de la Coh\'esion des Territoires et des Relations avec les Collectivit\'es Territoriales, co-direction avec M.-L. Breuill\'e (CESAER) \\
2016 - 2021 & \og{}\textit{Integrative mechanisms for addressing spatial justice and territorial inequalities in Europe}\fg{}, H2020 Project coordonn\'e par l'Universit\'e d'Aberystwyth (Pays de Galles), direction scientifique de l'\'equipe de recherche du CESAER \\
2016 - 2017 & \og{}\textit{Les d\'eterminants de la formation des niveaux de loyers}\fg{}, Projet de recherche financ\'e par le Minist\`ere du Logement, de l'Egalit\'e des Territoires et de la Ruralit\'e, co-direction avec M.-L. Breuill\'e (CESAER) et I.~Nappi-Choulet (ESSEC) \\
2015 - 2019 & \og{}\textit{Les b\'en\'efices de la densification urbaine dans la perspective du Grand Paris Express}\fg{}, contrat de recherche avec la Soci\'et\'e du Grand Paris, co-direction avec M.-L. Breuill\'e (CESAER) \\
2012 & \og{}\textit{R\'eorganisation des activit\'es productives et recomposition des territoires dans l'espace euro-m\'editerran\'een}\fg{}, BQR d\'ecern\'e par l'Universit\'e de Franche-Comt\'e \\
2009 & \og{}\textit{Nature et organisation des activit\'es industrielles et d'innovation, le r\^ole des \'economies d'agglom\'eration}\fg{}, BQR d\'ecern\'e par l'Universit\'e de Franche-Comt\'e \\
2006 & \og{}\textit{Cr\'eation d'un nouvel axe de recherche au sein du CRESE}\fg{}, BQR d\'ecern\'e par l'Universit\'e de Franche-Comt\'e \\

\end{longtable}

\smalltitre{Participation \`a des contrats de recherche}

\noindent \begin{longtable}{@{}p{\annee}p{\texte}@{}}
2026 - 2029 & \og{}\textit{Evolution des produits stup\'efiants saisis en France sur la longue dur\'ee}\fg{}, Programme Interminist\'eriel de Recherches Appliqu\'ees \`a La Lutte Antidrogue, dirig\'e par F.~Besancier (Laboratoire de Police Scientifique de Lyon) \\
2026 - 2030 & \og{}\textit{Market Access, Historical Development \& Inequalities}\fg{}, ANR 2025 Appel g\'en\'erique, dirig\'e par T.~Th\'evenin (THEMA, Universit\'e Bourgogne Europe) \\
2026 - 2031 & \og{}\textit{Live here or leave - (Un)Expected Impacts of Residential Location: Health, time and Financial}\fg{}, ANR 2025 Appel g\'en\'erique, dirig\'e par H.~Bouscasse (CESAER, INRAE) \\
2025 - 2029 & \og{}\textit{RUral Revival: A Life In The Inspirational Countryside}\fg{}, HORIZON Project, dirig\'e par G.~Lafert\'e (CESAER, INRAE) \\
2024 - 2029 & \og{}\textit{Land Market Regulations and Territorial Inequalities}\fg{}, ANR 2023 Appel g\'en\'erique, dirig\'e par G.~Chapelle (THEMA, CY Cergy Paris Universit\'e) \\
2022 - 2025 & \og{}\textit{Designing a Roadmap for Effective and Sustainable Strategies for Assessing and Adressing the Challenges of EU Agriculture to Navigate within a Safe and Just Operating Space}\fg{}, H2020 Project, dirig\'e par L.~V\'edrine pour le CESAER \\
2022 - 2023 & \og{}\textit{Impact \'economique des espaces verts \`a Qu\'ebec : Investigation sur la base du march\'e r\'esidentiel unifamilial}\fg{}, Projet de recherche financ\'e par le la Ville de Qu\'ebec, dirig\'e par J.~Dub\'e \\
2020- 2023 & \og{}\textit{D\'emat\'erialisation et acc\`{e}s \`a la justice}\fg{}, DEMAJUST 2020, dirig\'e par M.~Hilal pour le CESAER \\
2019	 - 2020 & \og{}\textit{Centralit\'es : comment les identifier et quels r\^oles dans les dynamiques locales et intercommunales ?}\fg{}, Projet de recherche financ\'e par le CGET, dirig\'e par M.~Hilal pour le CESAER \\
2018 - 2021 & \og{}\textit{PubPrivLand: Interactions public-priv\'e - causes \'economiques, sociales et environnementales des diff\'erents usages et modes de gestion}\fg{}, projet Int\'egr\'e  de l'Axe 2-ISITE-BFC, dirig\'e par F.~Raoul (Chrono-environnement, UFC) et N. Renahy (CESAER, INRAE) \\
2018 - 2022 & \og{}\textit{Low-Input Farming and Territories}\fg{}, H2020 Project, dirig\'e par S.~Legras et L.~V\'edrine pour le CESAER \\
2017 - 2018 & \og{}\textit{Parents Isol\'es Recherche Appartement : Test et Evaluation}\fg{}, Projet de recherche financ\'e par la Ville de Paris, dirig\'e par P.~Petit (ERUDITE, Universit\'e Paris-Est Marne-la-Vall\'ee) \\
2016 - 2018 & \og{}\textit{Discrimination dans l'acc\`{e}s au logement : un testing de couverture nationale}\fg{}, ANR 2015 Appel g\'en\'erique, dirig\'e par Y.~L'Horty (ERUDITE, Universit\'e Paris-Est Marne-la-Vall\'ee) \\
2015 - 2016 & \og{}\textit{Les d\'eterminants de la formation des niveaux de loyers}\fg{}, Projet de recherche financ\'e par le Minist\`{e}re du Logement, de l'Egalit\'e des Territoires et de la Ruralit\'e, dirig\'e par I.~Nappi-Choulet (ESSEC) \\
2015 - 2016 & \og{}\textit{Evaluation sociologique et statistique des actions innovantes pour la mobilit\'e et l'accompagnement des jeunes vers l'emploi}\fg{}, Fonds d'Exp\'erimentation pour la Jeunesse, dirig\'e par Y.~L'Horty (ERUDITE, Universit\'e Paris-Est Marne-la-Vall\'ee) \\
2015 		& \og{}\textit{Mod\'elisation des interactions spatiales entre les communes fran\c{c}aises via la coop\'eration intercommunale}\fg{}, PEPS MoMIS, dirig\'e par M.~Barth\'el\'emy (Institut de Physique Th\'eorique, CEA) \\
2014 - 2016 & \og{}\textit{Fiscalit\'e \'energ\'etique, territoires et f\'ed\'eralisme fiscal}\fg{}, Conseil Fran\c{c}ais de l'Energie, dirig\'e par E.~Taugourdeau (CES-ENS Cachan, PSE) \\
2012 - 2015 & \og{}\textit{Usage des sols : mod\`{e}les, dynamique et d\'ecisions}\fg{}, Programme Blanc 2011, Agence Nationale de la Recherche, dirig\'e par C.~Thomas-Agnan (GREMAQ, UMR-CNRS 5604, Universit\'e de Toulouse 1) \\
2011 - 2014 & \og{}\textit{Fiscal Decentralization and Public Private Partnerships, A Means of Improving Efficiency in the Delivery of Public Goods and Services?}\fg{}, Projet financ\'e par le FNSRS regroupant : IDHEAP/Unil (Suisse), Universit\'e de Fribourg (Suisse), Universit\'e de Franche-Comt\'e, Universit\'e Paris 1, Universit\'e de Lugano (Suisse), dirig\'e par L.~Athias (Universit\'e de Paris 1) et T.~Madi\`{e}s (Universit\'e de Fribourg, Suisse)\\
2010 - 2013 & \og{}\textit{Causality, Outliers and Nonlinearity in Spatiotemporal Models}\fg{}, Projet financ\'e par : Universidad de Zaragoza (Espagne), dirig\'e par J.~Mur (Universidad de Zaragoza, Espagne)\\
2009 - 2010 & \og{}\textit{Analysis Economico Espacial del Precio de la Vivienda en la Ciudad de Madrid}\fg{}, Projet financ\'e par : Universidad Aut\'onoma de Madrid (Espagne), dirig\'e par C.~Chasco (UAM, Espagne)\\
2007 - 2008 & \og{}\textit{L'innovation en Franche-Comt\'e - Pour une meilleure orientation des politiques publiques 2007-2013}\fg{}, Projet financ\'e par : Pr\'efecture de la R\'egion Franche-Comt\'e / SGAR, dirig\'e par F.~Picard (RECITS, UTBM), prestataire : MSHE Claude-Nicolas Ledoux \\
2006 - 2008 & \og{}\textit{Dynamiques r\'egionales et gouvernance au sein de l'Union Europ\'eenne Elargie }\fg{}, Contrat de Projet Etat-R\'egion, dirig\'e par C.~ Ertur et R.~Guillain (LEG, UMR CNRS 5118 Universit\'e de Bourgogne) \\
2005 - 2008 & \og{}\textit{Dynamiques r\'egionales, territoires urbains et modes de gouvernance au sein de l'Union Europ\'eenne \'elargie}\fg{}, Programme Jeunes Chercheuses et Jeunes Chercheurs 2005, Agence Nationale de la Recherche, dirig\'e par R.~Guillain (LEG, UMR-CNRS 5118, Universit\'e de Bourgogne) \\
2004 - 2006 & \og{}\textit{Les territoires urbains polycentriques en Europe. De l'organisation \'economique \`a la gouvernance}\fg{}, Plan Urbanisme Construction Architecture, Minist\'ere de l'\'equipement, des transports, de l'am\'enagement du territoire et de la mer, dirig\'e par C.~Baumont (LEG, UMR-CNRS 5118, Universit\'e de Bourgogne) \\
2003 - 2006 & \og{}\textit{La valeur \'economique des paysages dans les villes p\'eriurbaines}\fg{}, Minist\'ere de l'\'ecologie et du d\'eveloppement durable, Programme APR S3E 2002, dirig\'e par H.~Jayet (MEDEE, Universit\'e de Lille) \\
2003 - 2005 & \og{}\textit{Les mutations r\'egionales et urbaines dans les pays d'Europe Centrale et Orientale}\fg{}, Axe \og{}\textit{Analyse spatiale des disparit\'es \'economiques r\'egionales en Europe face \'e l'int\'egration des PECO}\fg{}, Contrat de Plan Etat-R\'egion, dirig\'e par C.~Ertur (LEG, UMR-CNRS 5118, Universit\'e de Bourgogne) \\
2003 		& \og{}\textit{Les politiques de d\'eveloppement r\'egional en Aquitaine}\fg{}, Conseil R\'egional d'Aquitaine, dirig\'e par J.~Le Cacheux (CATT, Universit\'e de Pau et des Pays de l'Adour) \\
2002-2006 	& \og{}\textit{Analyse spatiale des ph\'enom\`{e}nes \'economiques}\fg{}, Groupement de Recherches, CNRS (GDR 2345), dirig\'e par H.~Jayet (MEDEE, Universit\'e de Lille) \\
2002-2003 	& Center for Spatially Integrated Social Science, National Science Foundation BCS 9978058, dirig\'e par M.J.~Goodchild (University of California, Santa Barbara, Etats-Unis) \\
2001-2003 	& \og{}\textit{Services aux entreprises et nouvelles centralit\'es urbaines}\fg{}, Plan Urbanisme Construction Architecture, Minist\`ere de l'Equipement, dirig\'e par J.-M.~Huriot (LATEC, UMR-CNRS 5118, Universit\'e de Bourgogne) \\
2000-2002 	& \og{}\textit{Analyse \'economique des nouvelles formes de suburbanisation dans les grandes villes fran\'eaises}\fg{}, Aide \`{a} Projets Nouveaux, CNRS, dirig\'e par C.~Baumont (LATEC, UMR-CNRS 5118, Universit\'e de Bourgogne)
\end{longtable}

%%% Animations de la recherche %%%

\smalltitre{Animation de la recherche} \medskip

\noindent - Membre du Comit\'e Scientifique, \textit{24\textsuperscript{th} Spatial Econometrics and Statistics Workshop}, Paris, 17-18 juin 2026\\ 
- Membre du Comit\'e Scientifique, \textit{8\textsuperscript{th} Geography of Innovation Conference}, Budapest (Hongrie, 28-30 janvier 2026	 \\
- Membre du Comit\'e Scientifique, \textit{23\textsuperscript{rd} Spatial Econometrics and Statistics Workshop}, Saint-Etienne, 4-5 juin 2025\\ 
- Membre du Comit\'e Scientifique, \textit{22\textsuperscript{nd} Spatial Econometrics and Statistics Workshop}, Grenoble, 21-22 mai 2024\\ 
- Membre du Comit\'e Scientifique, \textit{14\textsuperscript{th} World Congress of the RSAI}, Kecskemet (Hongrie), 8-11 avril 2024\\
- Membre du Comit\'e Scientifique, \textit{7\textsuperscript{th} Geography of Innovation Conference}, Manchester (Royaume-Uni), 10-12 janvier 2024\\
- Membre du Comit\'e Scientifique, \textit{XVII\textsuperscript{th} World Conference of the Spatial Econometrics Association}, San Diego (Etats-Unis), 16-17 novembre 2023 \\
- Membre du Comit\'e Scientifique, \textit{21\textsuperscript{th} Spatial Econometrics and Statistics Workshop}, Dijon, 25-26 mai 2023\\ 
- Membre du Comit\'e Scientifique, \textit{11\textsuperscript{th} Seminar Jean Paelinck}, Cartagena (Espagne), 19-20 avril 2023\\ 
- Membre du Comit\'e Scientifique, \textit{20\textsuperscript{th} Spatial Econometrics and Statistics Workshop}, Lille, 19-20 mai 2022\\  
- Membre du Comit\'e Scientifique, \textit{6\textsuperscript{th} Geography of Innovation Conference}, Milan (Italie), 26-28 janvier 2022 \\
- Membre du Comit\'e Scientifique, \textit{19\textsuperscript{th} Spatial Econometrics and Statistics Workshop}, Nantes, 31 mai - 1er juin 2021\\
- Membre du Comit\'e Scientifique, \textit{13\textsuperscript{th} World Congress of the RSAI}, Marrakech (Maroc), 25-28 mai 2021\\
- Membre du Comit\'e Scientifique, \textit{5\textsuperscript{th} Geography of Innovation Conference}, Stavanger (Norv\`{e}ge), 31 janvier - 2 f\'evrier 2020	 \\
- Membre du Comit\'e Scientifique, \textit{59\textsuperscript{th} ERSA Congress}, Lyon, 27-30 ao\^ut 2019\\
- Membre du Comit\'e Scientifique, \textit{18\textsuperscript{th} Spatial Econometrics and Statistics Workshop}, Paris, 23-24 mai 2019\\
- Membre du Comit\'e Scientifique, \textit{5\textsuperscript{th} Global Conference in Economic Geography 2018}, Cologne (Allemagne), 24-28 juillet 2018\\
- Membre du Comit\'e Scientifique, \textit{55\textsuperscript{\`eme} colloque de l'Association de Science R\'egionale de Langue Fran\c{c}aise}, Caen, 4-6 juillet 2018\\
- Membre du Comit\'e Scientifique, \textit{XII\textsuperscript{th} World Conference of the Spatial Econometrics Association}, Vienne (Autriche), 11-12 juin 2018\\
- Organisation du \textit{17\textsuperscript{th} Spatial Econometrics and Statistics Workshop}, Dijon, 24-25 mai 2018\\
- Membre du Comit\'e Scientifique, \textit{4\textsuperscript{th} Geography of Innovation Conference}, Universitat de Barcelona (Espagne), 31 janvier - 2 f\'evrier 2018\\
- Membre du Comit\'e Scientifique, \textit{9\textsuperscript{th} Seminar Jean Paelinck}, Cartagena (Espagne), 6-7 octobre 2017\\
- Membre du Comit\'e Scientifique, \textit{11\textsuperscript{th} World Conference of the Spatial Econometrics Association}, Singapore Management University, 13-15 juin 2017\\
- Membre du Comit\'e Scientifique, \textit{16\textsuperscript{th} Spatial Econometrics and Statistics Workshop}, Avignon, 18-19 mai 2017\\
- Membre du Comit\'e Scientifique, \textit{3\textsuperscript{rd} Geography of Innovation Conference}, Sciences Po Toulouse (France), 28-30 janvier 2016\\
- Membre du Comit\'e Scientifique, \textit{15\textsuperscript{th} Spatial Econometrics and Statistics Workshop}, Orl\'eans, 26-27 mai 2016\\
- Membre du Comit\'e Scientifique, \textit{Workshop INFER Spatial Analysis in Energy, Environmental and Innovation Economics}, Universit\'e of Aachen (Allemagne), 10-11 d\'ecembre 2015\\
- Membre du Comit\'e Scientifique, \textit{8\textsuperscript{th} Seminar Jean Paelinck}, Dijon (France), 23-24 octobre 2015\\
- Membre du Comit\'e Scientifique, \textit{Workshop INFER Economic Geography, Agglomeration and Environment}, Universit\'e de Corse, 24-25 septembre 2015\\
- Membre du Comit\'e Scientifique, \textit{4\textsuperscript{th} Global Conference in Economic Geography 2015}, Oxford (Royaume-Uni), 19-23 ao\^ut 2015\\
- Membre du Comit\'e Scientifique, \textit{14\textsuperscript{th} Spatial Econometrics and Statistics Workshop}, Paris, 27-28 mai 2015\\
- Membre du Comit\'e Scientifique, \textit{7\textsuperscript{th} Seminar Jean Paelinck}, Zaragoza (Espagne), 20-21 novembre 2014\\
- Membre du Comit\'e Scientifique, \textit{11\textsuperscript{\`eme} Conference Annuelle Territoires, Espaces et Politiques Publiques}, Nantes, 25-26 septembre 2014\\
- Membre du Comit\'e Scientifique, \textit{13\textsuperscript{th} Spatial Econometrics and Statistics Workshop}, Toulon, 15-16 avril 2014\\
- Membre du Comit\'e Scientifique, \textit{10\textsuperscript{th} World Congress of the RSAI}, Bangkok (Thailande), 27-29 mai 2014\\
- Membre du Comit\'e Scientifique, \textit{2\textsuperscript{nd} Eurolio European Seminar Geography of Innovation}, Universit\'e d'Utrecht (Pays-Bas), 23-25 janvier 2014\\
- Membre du Comit\'e Scientifique, \textit{6\textsuperscript{th} Seminar Jean Paelinck}, Madrid (Espagne), 18-19 octobre \\
- Membre du Comit\'e Scientifique, \textit{12\textsuperscript{th} Spatial Econometrics and Statistics Workshop}, Orl\'eans, 17-19 juin 2013\\
- Membre du Comit\'e Scientifique, \textit{11\textsuperscript{th} Spatial Econometrics and Statistics Workshop}, Avignon, 15-16 novembre 2012\\
- Membre du Comit\'e Scientifique, \textit{5\textsuperscript{th} Seminar Jean Paelinck}, Coimbra (Portugal), 26-27 octobre 2012\\
- Membre du Comit\'e Scientifique, \textit{52\textsuperscript{nd} Congress of the European Regional Science Association}, University of Bratislava (Slovaquie), 21-25 ao\^ut 2012\\
- Organisatrice et pr\'esidente de la session \og~Econom\'etrie spatiale~\fg, \textit{XLIX\textsuperscript{e} Colloque de l'ASRDLF, Industrie, villes et r\'egions dans une \'economie mondialis\'ee}, Belfort, 9-11 juillet 2012\\
- Membre du Comit\'e Scientifique, \textit{Territoires, Emplois et Politiques Publiques}, Universit\'e de Caen Basse-Normandie, 14-15 juin 2012\\
- Membre du Comit\'e Scientifique, \textit{Eurolio European Seminar Geography of Innovation}, Universit\'e de Saint-Etienne, 26-28 janvier 2012\\
- Membre du Comit\'e Scientifique, \textit{Territoires, Emplois et Politiques Publiques}, Universit\'e de Metz, 23-24 juin \\
- Membre du Comit\'e Scientifique, \textit{V\textsuperscript{th} World Conference of the Spatial Econometrics Association}, Universit\'e de Toulouse I, 6-8 juillet 2011\\
- Membre du Comit\'e Scientifique, \textit{4\textsuperscript{th} Seminar Jean Paelinck}, Oviedo (Espagne), 22-23 octobre 2010\\
- Membre du Comit\'e Scientifique, \textit{Colloques joints ASRDLF and AISRe}, Aoste (Italie), 20-22 septembre 2010\\
- Membre du Comit\'e Scientifique, \textit{IV\textsuperscript{th} World Conference of the Spatial Econometrics Association}, Chicago (Etats-Unis), 9-12 juin 2010\\
- Membre du Comit\'e Scientifique, \textit{9\textsuperscript{th} Spatial Econometrics and Statistics Workshop}, Orl\'eans, 24-25 juin 2010\\
- Organisation du \textit{8\textsuperscript{th} Spatial Econometric and Statistics Workshop}, Besan\c{c}on, 2-3 juin 2009\\
- Membre du Comit\'e Scientifique, \textit{3\textsuperscript{rd} World Conference of the Spatial Econometrics Association}, Barcelone (Espagne), 8-10 juillet 2009\\
- Membre du Comit\'e Scientifique, \textit{14\textsuperscript{th} International Conference on Computing in Economics and Finance}, Paris, 26-28 juin 2009\\
- Membre du Comit\'e Scientifique, \textit{7\textsuperscript{th} Spatial Econometric and Statistics Workshop}, Paris, 9-10 juin 2009\\
- Organisatrice des sessions d'\'econom\'etrie et de statistique spatiales, \textit{54\textsuperscript{th} North American Meetings of the Regional Science Association International}, Savannah (Etats-Unis), 7-11 novembre 2007(avec A.~P\'aez, S.~Dall'erba et R.~Buliung) \\
- Pr\'esidente et rapportrice de la session EPAINOS \og~Agglomeration~\fg, \textit{47\textsuperscript{th} Congress of the European Regional Science Association}, Paris, 29 ao\^ut~~2 septembre 2007\\
- Membre du Comit\'e Scientifique, \textit{6\textsuperscript{th} Spatial Econometric and Statistics Workshop}, Dijon, 4-5 juin 2007\\
- Membre du Comit\'e Scientifique, \textit{ADRES International Conference, Networks of Innovation and Spatial Analysis of Knowledge Diffusion}, Saint-Etienne, 14-15 septembre 2006\\
- Membre du Comit\'e Scientifique, \textit{5\textsuperscript{th} Spatial Econometric and Statistics Workshop}, Grenoble, 14 juin 2006 \\

\smalltitre{Transfert de la recherche} \medskip

\noindent 
- \textbf{Cabinet CASALONGA}. Communication : \og{}Les d\'eterminants du foncier viticole et sa valorisation\fg{} (20~mars 2025) \\
- \textbf{NEOTOA Ile et Vilaine}. Communication : \og{}Depuis la crise de la Covid-19, le destinataire n'habite plus \`a  l'adresse indiqu\'ee\fg{} (17~octobre 2024) \\
- \textbf{Union Sociale pour l'Habitat}. Communication : \og{}Depuis la crise de la Covid-19, le destinataire n'habite plus \`a  l'adresse indiqu\'ee\fg{} (28~mai 2024) \\
- \textbf{Institut des Hautes Etudes pour l'Action dans le Logement}. Communication : \og{}C'est quoi le bail ? Observer, comparer, comprendre et r\'eguler les loyers en France\fg{} (18~janvier 2024) \\
- \textbf{Seminaire Economie G\'eographique INSEE}. Communications : \og{}Des arbitrages micro\'economiques \'a l'\'eco-gentrification : Enjeux environnementaux de la mobilit\'e r\'esidentielle\fg{}  et \og{}Eclairer les politiques publiques : Quels d\'efis pour les \'etudes territoriales et l'analyse spatiale\fg{} (29~septembre 2023) \\
- \textbf{Installation du groupe de travail permanent sur les meubl\'es de tourisme}. Communication : \og{}Pr\'esentation d'une \'etude sur le lien entre meubl\'es de tourisme et loyers\fg{} (23~mars 2021) \\
- \textbf{Conf\'erence des citoyens sur les locations de meubl\'es de tourisme}. Communication : \og{}L'effet inflationniste des locations intensives sur le march\'e locatif\fg{} (23~janvier 2021) \\
- \textbf{Rencontres de l'Observation des Loyers}. Communication : \og{}Des donn\'ees pour la recherche\fg{} (8~f\'evrier 2018) \\
- \textbf{Inter Rh\^one}. Communication : \og{}Estimation des \'elasticit\'es-prix / quantit\'es pour les vins de C\^otes du Rh\'ene : march\'e du vrac et grande distribution\fg{} (12~f\'evrier 2015) \\
- \textbf{Syndicat des Transports en Ile-de-France}. Communication : \og{}Les territoires d'emplois en Ile-de-France : persistance et mutations des polarisations\fg{} (8~juillet 2006) \\
- \textbf{CROCIS, Chambre de Commerce et d'Industrie de Paris}. Communication : \og{}Syst\`{e}mes d'Information G\'eographique et localisation des emplois en Ile-de-France\fg{} (20 avril 2005)

%%% Direction et jury de th\'eses de doctorat %%%

\titre{Encadrement doctoral et d'HDR \\ participation \`{a} des jurys de th\`eses de doctorat et d'HDR}

\smalltitre{Encadrement de th\`{e}ses de doctorat}
	
\noindent \begin{longtable}{@{}p{\annee}p{\texte}@{}}
2016 - 2020 & Oudah~YOBOM, Th\`{e}se de Doctorat (CESAER, UMR1041, Universit\'e de Bourgogne Franche-Comt\'e) \\
\end{longtable}

\smalltitre{Co-encadrement de th\`{e}ses de doctorat}
	
\noindent \begin{longtable}{@{}p{\annee}p{\texte}@{}}
2024 -		& Simon~DELAGE Th\`{e}se de Doctorat CIFRE co-encadr\'ee avec M.-L.~Breuill\'e (CESAER, UMR 1041, INRAE) et A.~Hamel (Geoptis) \\
2024 -		& Fares~BEN YOUSSEF Th\`{e}se de Doctorat co-encadr\'ee avec M.~Kriaa (Universit\'e Tunis, Tunisie) \\
2023 -		& Ivan~MARIN Th\`{e}se de Doctorat en co-tutelle avec l'Universit\'e Laval, co-encadr\'ee avec J.~Dub\'e (Universit\'e Laval) \\
2023 -		& Louise~PHUNG, Th\`{e}se de Doctorat co-encadr\'ee avec M.-L.~Breuill\'e (CESAER, UMR 1041, INRAE) \\
2023 -		& Francois~LAFONT, Th\`{e}se de Doctorat co-encadr\'ee avec M.-L.~Breuill\'e (CESAER, UMR 1041, INRAE) \\
2022 - 2025	& Alexandra~VERLHIAC, Th\`{e}se de Doctorat co-encadr\'ee avec M.-L.~Breuill\'e (CESAER, UMR 1041, INRAE) et B.~Castillo (Meilleurs Agents) \\
2020 - 2025	& Martin~REGNAUD, Th\`{e}se de Doctorat CIFRE co-encadr\'ee avec M.-L.~Breuill\'e (CESAER, UMR 1041, INRAE) et P.~Vidal (Meilleurs Agents) \\
2020 - 2023	& Capucine~CHAPEL, Th\`{e}se de Doctorat co-encadr\'ee avec M. Hilal (CESAER, UMR 1041, INRAE) \\
2020 - 2023	& Yacine~ALLAM, Th\`{e}se de Doctorat co-encadr\'ee avec M.-L.~Breuill\'e (CESAER, UMR 1041, INRAE) \\
2019 - 2023	& R\'emi~LEI, Th\`{e}se de Doctorat co-encadr\'ee avec J.-S. Ay (CESAER, UMR 1041, INRAE) \\
2017 - 2022	& Rawaa~LAAJIMI, Th\`{e}se de Doctorat en co-tutelle avec l'Universit\'e de Sousse (Tunisie), co-encadrement avec S. Benammou (D\'epartement de Gestion et des M\'ethodes Quantitatives, Universit\'e de Sousse, Tunisie) \\
2016 - 2019 & Alexandre~FLAGE, Th\`{e}se de Doctorat co-encadr\'ee avec F. Cochard (CRESE, EA 3190, Universit\'e de Franche-Comt\'e) \\
2014 - 2018 & Youba~N'DIAYE, Th\`{e}se de Doctorat co-encadr\'ee avec M.-L.~Breuill\'e (CESAER, UMR 1041, INRA) \\
2013 - 2016 & Jean-Christian~TISSERAND, Th\`{e}se de Doctorat co-encadr\'ee avec Y.~Gabuthy (CRESE, EA 3190, Universit\'e de Franche-Comt\'e) \\
2008 - 2014 & Benjamin~LAURENT, Th\`{e}se de Doctorat co-encadr\'ee avec R.~Guillain (LEG, UMR-CNRS 5118, Universit\'e de Bourgogne) \\
2005 - 2013 & La\"{i}sa~RO'I, Th\`{e}se de Doctorat co-encadr\'ee avec M.-A.~S\'en\'egas (GREThA, UMR-CNRS 5113, Universit\'e de Bordeaux IV)
\end{longtable}

%%% Jurys de th\'ese %%%

\smalltitre{Participation \`a des jurys de th\`{e}se de doctorat} \medskip
	
\noindent 
- Olivier~DENAGISCARDE, Universit\'e Paris 1 Panth\'eon Sorbonne, suffrageante, 2025 \\
- Valentin~COCCO, Universit\'e Paris Saclay, rapporteure de la th\`{e}se, 2025 \\
- Etienne~DE L'ESTOILE, Universit\'e Paris 1, suffrageante, 2025 \\
- Benjamin~COTTREAU, Universit\'e de Lyon 2, rapporteure de la th\`{e}se, 2025 \\
- Martin~FAULQUES, Universit\'e de Normandie, rapporteure de la th\`{e}se, 2024 \\
- Mehdi~BERRADA, Universit\'e de Montpellier, rapporteure de la th\`{e}se, 2024 \\
- Calum~ROBERTSON, Universit\'e La Rochelle, rapporteure de la th\`{e}se, 2024 \\
- Nicolas~JUSTE, Universit\'e Lille, rapporteure de la th\`{e}se, 2024 \\
- Anran~MAO, Universit\'e Paris Saclay, rapporteure de la th\`{e}se, 2023 \\
- Florian~FOUQUET, Nantes Universit\'e, rapporteure de la th\`{e}se, 2023 \\
- David~DADAKPETE, Universit\'e de Bourgogne Franche-Comt\'e, 2023 \\
- Pedro Herrera~CATALAN, Universitat Aut\`{o}noma de Madrid (Espagne), rapporteure de la th\`{e}se, 2021 \\
- Najla~KAMERGI, Universit\'e de Toulon, pr\'esidente du jury, 2021 \\
- Juan Andr\'es Piedra~PENA, Universitat Aut\`{o}noma de Barcelona (Espagne), 2021 \\ 
- Darya~KELES, Universit\'e de Lorraine, rapporteure de la th\`{e}se, 2021 \\ 
- Yuheng~LING, Universit\'e de Corse, rapporteur de la th\`{e}se, 2020 \\ 
- Oguzhan~AKGUN, Universit\'e Paris 2, rapporteure de la th\`{e}se, 2019 \\ 
- Benjamin~EGRON, Universit\'e Paris Ouest, rapporteure de la th\`{e}se, 2019 \\ 
- Ouassim~MANOUT, Universit\'e Lyon III, rapporteure de la th\`{e}se, 2019 \\ 
- Emad~Aldeen~DARWICH, Universit\'e Lille III, rapporteur de la th\`{e}se, 2017 \\ 
- Giulia~CARRA, Universit\'e Paris-Saclay, 2017 \\ 
- Laurent~BERGE, Universit\'e de Bordeaux, rapporteure de la th\`{e}se, 2015 \\ 
- Laeticia~TUFFERY, Universit\'e Evry Val d'Essonne, rapporteure de la th\`{e}se, 2015 \\ 
- Sileymane~BA, Universit\'e de Bourgogne, 2015 \\ 
- Aligui~TIENTAO, Universit\'e de Bourgogne, pr\'esidente du jury, 2015 \\ 
- Moussa~FALL, Universit\'e d'Aix-Marseille, rapporteure de la th\`{e}se, 2015 \\ 
- Abdourahmanne~DIALLO, Universit\'e d'Aix-Marseille, rapporteure de la th\`{e}se, 2014 \\ 
- Hermann Pythagore Pierre~DONFOUET, Universit\'e de Rennes 1, rapporteure de la th\`{e}se, 2013 \\ 
- Marc~BRUNETTO, Universit\'e de Toulon, rapporteure de la th\`{e}se, 2013 \\ 
- Walliya~PREMCHIT, Universit\'e de Strasbourg, rapporteure de la th\`{e}se, 2013 \\ 
- Florent~SARI, Universit\'e Paris-Est, rapporteure de la th\`{e}se, 2011 \\ 
- Jean-Sauveur~AY, Universit\'e de Bourgogne, pr\'esidente du jury, 2011 \\ 
- Nejla~BEN ARFA, AgroParisTech, rapporteure de la th\`{e}se, 2011 \\
- Mohamed~AMARA, Universit\'e de Paris I et Universit\'e de Tunis, Tunisie, rapporteure de la th\`{e}se, 2010 \\
- Angela~PARENTI, Universit\`a di Pisa (Italie), pr\'esidente du jury, 2010 \\ 
- Octavio~ESCOBAR, Universit\'e Paris-Dauphine, rapporteure de la th\`{e}se, 2009 \\
- Alexandra~SCHAFFAR, Universit\'e de la R\'eunion, rapporteure de la th\`{e}se, 2009 \\ 
- Yacine~LARBI, Universit\'e Jean Monnet de Saint-Etienne et Institut National de la Planification et de la Statistique d'Alger, Alg\'erie, rapporteure de la th\`{e}se, 2009 \\ 
- Benjamin~NDONG, Universit\'e de Franche-Comt\'e, pr\'esidente du jury, 2007 \\
- Wilfried~KOCH, Universit\'e de Bourgogne, rapporteure de la th\`{e}se, 2007 \\
- Marie~LEBRETON, Universit\'e de la M\'editerran\'ee, 2006 \\

\smalltitre{Encadrement d'HDR} \medskip

\noindent - Nicolas DEBARSY, CNRS, LEM Universit\'e Lille, 2023 \\
\noindent - Raja CHAKIR, INRAE, AgroParisTech, 2013 \\

\smalltitre{Participation \`a des jurys d'HDR} \medskip

\noindent 
- Pascal RICORDEL, Universit\'e du Havre Normandie, rapporteure, 2026 \\
- Mariona SEGU, CY University, rapporteure, 2025 \\
- Ad\'ela\"ide FADHUILE, Universit\'e Grenoble-Alpes, pr\'esidente, 2024 \\
- Sylvain CHAREYRON, Universit\'e Paris-Est, rapporteure, 2023 \\
- Jeanne DACHARY-BERNARD, Universit\'e de Bourgogne Franche-Comt\'e, rapporteure, 2021 \\
- Elsa MARTIN, Universit\'e de Montpellier, 2017 \\
- Eric MARCON, Universit\'e de Guyane, rapporteure, 2016 \\
- Di\'ego LEGROS, Universit\'e d'Orl\'eans, rapporteure, 2015 \\
- Alexandra SCHAFFAR, Universit\'e de Toulon, rapporteure, 2014 \\
- Corinne AUTANT-BERNARD, Universit\'e de Saint-Etienne, pr\'esidente du jury, 2010 \\

\newpage

%%% Colloques %%%

\titre{Colloques et conf\'erences}

\noindent Seuls sont list\'es les colloques et conf\'erences auxquels j'ai particip\'e et pr\'esent\'e un article. \bigskip %%% Conf\'erences invit\'es %%%


\smalltitre{Conf\'erences invit\'ees} \medskip

\noindent - \textit{22\textsuperscript{th} Workshop Local Power, Local Growth, Taxes, Governance, and Business Synergy}, Paris, 3 d\'ecembre 2025\\
- \textit{22\textsuperscript{th} Spatial Econometrics Workshop}, Grenoble, 23-24 mai 2024, \textbf{keynote speaker}\\
- \textit{European Regional Science Association 62\textsuperscript{th} Congress 2023}, Alicante, Espagne, 28 ao\^ut-01 septembre 2023, \textbf{keynote speaker} \\
- \textit{Stochastic Geometry Days 2023}, RT GeoSto, Dijon, 12-16 septembre 2023, \textbf{keynote speaker} \\
- \textit{ESRC Workshop Recent Developments in Spatial Econometrics}, Royal Holloway University of London, 9-10 septembre 2022, \textbf{keynote speaker} \\
- \textit{XVI World Conference of Spatial Econometrics Association}, Universit\'e de Varsovie, 23-24 juin 2022, \textbf{keynote speaker} \\
- \textit{3\textsuperscript{rd} Wine and Spirits Economics Workshop}, Burgundy School of Business, 9 juin 2022 \\
- \textit{3\textsuperscript{rd} International Conference on Urban and Regional Economics}, Universit\'e de Katowice, Pologne, 12-13 juin 2015\\
- \textit{11\textsuperscript{th} Spatial Econometrics Workshop}, INRA Ecod\'eveloppement Avignon, 15-16 novembre 2012, \textbf{keynote speaker}\\
- \textit{Scale, location and spatial interactions in the economic analysis of multi-functional natural resources: Lessons for forestry}, Laboratoire d'Economie Foresti\`{e}re, AgroParisTech, Nancy, 16-17 septembre 2010\\
- \textit{2010 Barcelona Workshop on Regional and Urban Economics}, AQR (Grup d'An\'alisi Quantitativa Regional), University of Barcelona (Espagne), 9-10 septembre 2010\\
- \textit{Third Seminar of Spatial Econometrics JEAN PAELINCK}, Universidad Polit\'ecnica de Cartagena (Espagne), 10-11 octobre 2008\\
- \textit {Second Seminar of Spatial Econometrics JEAN PAELINCK}, Universit\'e de Saragosse (Espagne), 27-28 octobre 2006\\
- \textit{ADRES International Conference, Networks of Innovation and Spatial Analysis of Knowledge Diffusion}, CREUSET, Universit\'e de Saint-Etienne, 14-15 septembre 2006\\
- \textit{4\textsuperscript{th} Spatial Econometrics Workshop}, GREMAQ (UMR-CNRS 5604), Universit\'e de Toulouse I, 27 juin 2005\\
- \textit{First Seminar of Spatial Econometrics JEAN PAELINCK}, Universit\'e de Saragosse (Espagne), 22-23 octobre 2004\\
- \textit{WIDER Project Meeting on Spatial Inequality in Development}, Helsinki (Finlande), 29-30 mai 2003\\
- \textit{2\textsuperscript{nd} Spatial Econometrics Workshop}, LATEC (UMR-CNRS 5118), Universit\'e de Bourgogne, 13 mai 2003\\
- \textit{Cornell/LSE/WIDER Spatial Inequality and Development Conference}, London School of Economics CEP, Londres (Royaume-Uni), 27-30 juin 2002\\
- \textit{Spatial Econometrics Workshop}, GREMAQ (UMR-CNRS 5604), L.S.P. (UMR-CNRS 5583), Universit\'e de Toulouse I et III, 14 juin 2002 \\

\newpage

%%% Colloques nationaux et nternationaux %%%

\smalltitre{Colloques nationaux et internationaux} \medskip

\noindent - \textit{20\textsuperscript{th} Spatial Econometrics and Statistics Workshop}, Lille, 19-20 mai 2022\\  
- \textit{2021 Annual Conference of the International Association for Applied Econometrics}, Webconference, 22-25 juin 2021\\
- \textit{New Techniques and Technologies for Statistics 2021}, European Commission, 9-11 mars 2021\\
- \textit{60\textsuperscript{th} Annual Meeting of the European Regional Science Association}, Webconference, 25-27 ao\^ut 2020\\
- \textit{59\textsuperscript{th} Annual Meeting of the European Regional Science Association}, Lyon, 27-30 ao\^ut 2019\\
- \textit{68\textsuperscript{th} Annual Meeting of the French Economic Association}, Orl\'eans, 17-19 juin 2019\\
- \textit{18\textsuperscript{th} Spatial Econometrics and Statistics Workshop}, Paris, 23-24 mai 2019\\
- \textit{8\textsuperscript{th} Seminar Jean Paelinck}, Dijon, 23-24 octobre 2015\\
- \textit{14\textsuperscript{th} Spatial Econometrics and Statistics Workshop}, Paris, 27-28 mai 2015\\
- \textit{2015 Workshop on Real Estate Economics and Finance}, EBS Wiesbaden (Allemagne), 7 mars 2015\\
- \textit{11\textsuperscript{eme} Conf\'erence Annuelle Territoires, Espaces et Politiques Publiques}, Universit\'e de Nantes, 25-26 septembre 2014\\	
- \textit{8\textsuperscript{th} World Conference of the Spatial Econometrics Association}, Z\"{u}rich (Suisse), 11-13 juin 2014\\
- \textit{13\textsuperscript{th} Spatial Econometrics and Statistics Workshop}, Toulon (France), 15-16 avril 2014\\
- textit{2\textsuperscript{nd} Eurolio European Seminar Geography of Innovation}, Universit\'e d'Utrecht (Pays-Bas), 23-25 janvier 2014\\
- \textit{Sixth Seminar of Spatial Econometrics JEAN PAELINCK}, Universidad Aut\'onoma de Madrid (Espagne), 17-18 octobre 2013\\
- \textit{7\textsuperscript{th} World Conference of the Spatial Econometrics Association}, Washington (Etats-Unis), 10-12 juillet 2013\\
- \textit{9\textsuperscript{th} Spatial Econometrics and Statistics Workshop}, Orl\'eans, 24-25 juin 2010\\
- \textit{56\textsuperscript{th} North American Meetings of the Regional Science Association International}, San Francisco (Etats-Unis), 18-21 novembre 2009\\
- \textit{49\textsuperscript{th} European Congress of the Regional Science Association International}, L\'od\'z (Pologne), 25-29 ao\^ut2009 \\
- \textit{48\textsuperscript{th} Annual Meeting of the Western Regional Science Association}, Napa Valley (Etats-Unis), 22-25 f\'evrier 2009\\
- \textit{7\textsuperscript{th} Spatial Econometrics and Statistics Workshop}, Paris, 9-10 juin 2008\\
- \textit{54\textsuperscript{th} North American Meetings of the Regional Science Association International}, Savannah (Etats-Unis), 7-11 novembre 2007\\
- \textit{47\textsuperscript{th} Congress of the European Regional Science Association}, Paris, 29 ao\^ut - 2 septembre 2007\\
- \textit{The First World Conference of the Spatial Econometrics Association}, Fitzwilliam College, Cambridge (Royaume-Uni), 11-14 juillet 2007\\
- \textit{53\textsuperscript{rd} North American Meetings of the Regional Science Association International}, Toronto (Canada), 16-18 novembre 2006\\
- \textit{45\textsuperscript{th} Congress of the European Regional Science Association}, Vrije Universiteit, Amsterdam (Pays-Bas), 23-27 ao\^ut 2005\\
- \textit{Spatial Econometrics Workshop}, Kiel (Allemagne), 8-9 avril 2005\\
- \textit{51\textsuperscript{st} North American Meetings of the Regional Science Association International}, Seattle (Etats-Unis), 11-13 novembre 2004\\
- \textit{XL\ieme{} Colloque de l'ASRDLF, Convergence et Disparit\'es R\'egionales au sein de l'Espace Europ\'een}, Bruxelles (Belgique), 1-3 septembre 2004\\
- \textit{4\textsuperscript{th} Congress on Proximity Economics : Proximity, Networks and Co-ordination}, Universit\'e de la M\'editerran\'ee, 17-18 juin 2004\\
- \textit{3\textsuperscript{rd} Spatial Econometrics Workshop}, BETA (UMR-CNRS 7522), Universit\'e Louis Pasteur-Strasbourg I, 9 juin 2004\\
- \textit{50\textsuperscript{th} North American Meetings of the Regional Science Association International}, Philadelphie (Etats-Unis), 20-22 novembre 2003\\
- \textit{XXXIX\ieme{} Colloque de l'ASRDLF, Concentration et S\'egr\'egation ; Dynamiques et Inscriptions Territoriales}, Lyon, 1-3 septembre 2003\\
- \textit{8\textsuperscript{th} International Conference of the European Union Studies Association}, Nashville (Etats-Unis), 27-29 mars 2003\\
- \textit{49\textsuperscript{th} North American Meetings of the Regional Science Association International}, Porto Rico (Etats-Unis), 14-16 novembre 2002\\
- \textit{Colloque des Jeunes Econom\`{e}tres}, Mont-Saint-Odile, 2-4 mai 2002\\
- \textit{48\textsuperscript{th} North American Meetings of the Regional Science Association International}, Charleston (Etats-Unis), 15-17 novembre 2001\\
- \textit{10\ieme{} Journ\'ees du SESAME}, Dijon, 14-15 septembre 2000\\
- \textit{40\textsuperscript{th} Congress of the European Regional Science Association}, Barcelone (Espagne), 30 ao\^ut - 2 septembre 2000\\
- \textit{6\textsuperscript{th} World Congress of the Regional Science Association International}, Lugano (Suisse), 16-20 mai 2000\\
- \textit{46\textsuperscript{th} North American Meetings of the Regional Science Association International}, Montr\'eal (Canada), 11-14 novembre 1999\\
- \textit{XLVIII\ieme{} Congr\`{e}s Annuel de l'AFSE}, Paris, 24-25 septembre 1999\\
- \textit{9\ieme{} Journ\'ees du SESAME}, Saint-Etienne, 8-10 septembre 1999\\

\newpage

%%% Activit\'es d'enseignement %%%

\bigtitre{Activit\'es d'enseignement}

%%% Agrosup Dijon %%%
	
\titre{l'Institut Agro Dijon}

\noindent \begin{longtable}{@{}p{\annee}p{\texte}@{}}	
2024 - 		&  \textbf{Macro\'economie} (CM et TD, Convention avec l'Institut National de Formation des Personnels du Minist\`ere de l'Agriculture) \\
			&  \textbf{Acteurs et m\'ecanismes \'economiques} (CM, Formation agronomie par apprentissage, 1\iere{}~ann\'ee) \\
			& \textbf{Statistiques} (TD, tronc commun, 1\iere{}~ann\'ee en formation initiale) \\	
2024 -2025	&  \textbf{Entreprises agroalimentaires et pouvoirs publics} (CM, Formation agro-alimentaire par apprentissage, 2\ieme{}~ann\'ee) \\	
2022 - 		&  \textbf{Statistiques spatiales} (CM, dominante Data et Num\'erique pour l'Agriculture et l'Alimentation, 3\ieme{}~ann\'ee) \\
2019 -		& \textbf{Les enjeux socio-\'economiques des syst\`emes agri-alimentaires} (CM, tronc commun, 1\iere{}~ann\'ee en formation initiale) \\
2019 - 2024	& \textbf{Politique alimentaire} (Formation continue) \\
2019 - 2024	& \textbf{Politique de la concurrence} (CM, dominante Strat\'egies et Organisation des Fili\`eres et Entreprises Agricoles et Agroalimentaires, 3\ieme{}~ann\'ee) \\
2018 -		& \textbf{Diagnostics de territoires} (CM et TD, sp\'ecialit\'e agronomie, dominante AGIR sur les territoires, 3\ieme{}~ann\'ee) \\
2018 - 2021 & \textbf{Micro\'econom\'etrie d'\'evaluation des politiques publiques} (CM, sp\'ecialit\'e agronomie, dominante AGIR sur les territoires, 3\ieme{}~ann\'ee) \\
2018 - 2020	& \textbf{Statistiques} (TD, tronc commun, 2\ieme{}~ann\'ee en formation initiale) \\
2016 - 	& \textbf{Economie et politiques agricoles et d'environnement} (Formation continue) \\
2015 -	& \textbf{Statistiques avanc\'ees} (CM et TD, sp\'ecialit\'e agronomie, dominante AGIR sur les territoires, 3\ieme{}~ann\'ee) \\
2015 -2024	& \textbf{Politiques agricoles et de l'environnement} (CM et TD, sp\'ecialit\'e agronomie, 2\ieme{}~ann\'ee en formation initiale) \\
2015 - 2024	& \textbf{Acteurs et m\'ecanismes \'economiques} (CM et TD, sp\'ecialit\'e agronomie, 1\iere{}~ann\'ee en formation initiale) \\
2015 -		& \textbf{Acteurs et m\'ecanismes \'economiques} (CM, tronc commun, 1\iere{}~ann\'ee en formation initiale) \\
2016		& \textbf{Module politiques publiques laiti\`{e}res} (MS MIP) \\	
2015		& \textbf{Encadrement du module fili\`{e}re lait-fromage} (TD, sp\'ecialit\'e agronomie, 2\ieme{}~ann\'ee en formation initiale) \\
\end{longtable}

\newpage

\titre{Universit\'e de Marie et Louis Pasteur, Besan\c{c}on}

\smalltitre{U.F.R. Sciences Juridiques, Economiques, Politiques et de Gestion}

\noindent \begin{longtable}{@{}p{\annee}p{\texte}@{}}
2024 -	& \textbf{Statistique et \'econom\'etrie sur donn\'ees spatiales} (CM, Master Economie-Gestion, sp\'ecialit\'e Quantification et Analyse Economique, 2\ieme{}~ann\'ee) \\
2015 - 2017	& \textbf{Econom\'etrie appliqu\'ee} (CM, Master Economie-Gestion, sp\'ecialit\'e Charg\'e d'Etudes Economiques, 1\iere{}~ann\'ee) \\
2014 - 2015	& \textbf{Economie industrielle} (CM, Licence Economie-Gestion, 3\ieme{}~ann\'ee) \\
			& \textbf{Concurrence imparfaite} (CM, Licence Economie-Gestion, 2\ieme{}~ann\'ee) \\
2013 - 2015	& \textbf{Evaluation des politiques publiques} (CM, Master Economie-Gestion, sp\'ecialit\'e Charg\'e d'Etudes Economiques, 2\ieme{}~ann\'ee) \\
			& \textbf{Economie publique locale} (CM, Master Economie-Gestion, sp\'ecialit\'e Charg\'e d'Etudes Economiques, 2\ieme{}~ann\'ee) \\
			& \textbf{Initiation \`a la recherche} (CM, Licence Economie-Gestion, parcours d'excellence, 3\ieme{}~ann\'ee) \\
2012 - 2015	& \textbf{Econom\'etrie approfondie} (CM, Master Economie-Gestion, sp\'ecialit\'e Charg\'e d'Etudes Economiques, 1\iere{}~ann\'ee) \\
			& \textbf{Econom\'etrie} (CM, Master Economie-Gestion, sp\'ecialit\'es Charg\'e d'Etudes Economiques, Banque et E-achat et March\'es, 1\iere{}~ann\'ee) \\
2012 - 2014	& \textbf{Politiques \'economiques compar\'ees} (CM, Licence Economie-Gestion, 3\ieme{}~ann\'ee) \\
2012 - 2013	& \textbf{Evaluation des politiques publiques} (CM, cours doctoral) \\
			& \textbf{Remise \`a niveau en macro\'economie} (CM, Licence Economie-Gestion, 3\ieme{}~ann\'ee) \\
			& \textbf{Institutions et r\'egulation} (CM, Master Economie-Gestion, sp\'ecialit\'es Charg\'e d'Etudes Economiques, Banque et E-achat et March\'es, 1\iere{}~ann\'ee) \\
			& \textbf{Micro\'econom\'etrie} (CM, Master Economie-Gestion, sp\'ecialit\'e Charg\'e d'Etudes Economiques, 2\ieme{}~ann\'ee) \\
			& \textbf{Politique de la concurrence} (CM, Master Economie-Gestion, sp\'ecialit\'e Charg\'e d'Etudes Economiques, 2\ieme{}~ann\'ee) \\
2011 - 2012	& \textbf{Econom\'etrie appliqu\'ee} (CM et TD, Master Economie de la Firme et des March\'es, sp\'ecialit\'e Expertise Economique, 1\iere{}~ann\'ee, cours commun avec C. Refait-Alexandre) \\
			& \textbf{Probl\`{e}mes \'economiques contemporains} (TD, Licence Administration Economique et Sociale, 1\iere{}~ann\'ee) \\	
2009 - 2010	& \textbf{S\'eminaires} (CM, Master Economie de la Firme et des March\'es, sp\'ecialit\'e Expertise Economique, 2\ieme{}~ann\'ee, cours commun avec P.-H.~Morand) \\
			& \textbf{Commerce international} (CM, Licence Economie et Gestion, 2\ieme{}~ann\'ee) \\
			& \textbf{Renforcement en \'economie} (CM, Licence Economie et Gestion, 1\iere{}~ann\'ee) \\
2008 - 2012	& \textbf{Micro\'econom\'etrie} (CM, Master Economie de la Firme et des March\'es, sp\'ecialit\'e Expertise Economique, 2\ieme{}~ann\'ee) \\
			& \textbf{Econom\'etrie} (CM et TD, Master Economie de la Firme et des March\'es, sp\'ecialit\'e Expertise Economique et E-achat et March\'es, 1\iere{}~ann\'ee) \\
			& \textbf{Economie spatiale} (CM et TD, Licence Administration Economique et Sociale, 3\ieme{}~ann\'ee) \\ 
2008 - 2009	& \textbf{M\'ethodologie de la recherche} (CM, Master Economie de la Firme et des March\'es, sp\'ecialit\'e Expertise Economique, 2\ieme{}~ann\'ee, cours commun avec P.-H.~Morand) \\
		 	& \textbf{Econom\'etrie appliqu\'ee} (CM et TD, Master Economie de la Firme et des March\'es, sp\'ecialit\'e Expertise Economique, 1\iere{}~ann\'ee) \\
			& \textbf{Data mining} (CM et TD, Master Economie de la Firme et des March\'es, sp\'ecialit\'e Expertise Economique et E-achat et March\'es, 1\iere{}~ann\'ee) \\
2006 - 2008	& \textbf{Micro\'econom\'etrie appliqu\'ee} (CM, Master Sciences Economiques, sp\'ecialit\'e recherche Micro\'economie Appliqu\'ee, 2\ieme{}~ann\'ee) \\
			& \textbf{Mod\'elisation \'economique et informatique} (TD, Master Sciences Economiques, 1\iere{}~ann\'ee) \\
2006 - 2007	& \textbf{Macro\'economie} (CM, Licence Droit-Economie-Gestion, mention Economie et Gestion, 2\ieme{}~ann\'ee) \\
\end{longtable}

%%% IUT de Belfort-Montb\'eliard %%%

\smalltitre{I.U.T. de Belfort-Montb\'eliard}

\noindent \begin{longtable}{@{}p{\annee}p{\texte}@{}}
2008 - 2009 & \textbf{Data mining} (CM et TD, D\'epartement Techniques de Commercialisation, Licence professionnelle Technologies de l'Information et de la Communication Appliqu\'ees au Marketing et au Commerce) \\
2006 - 2008 & \textbf{Probl\'ematiques \'economiques appliqu\'ees} (CM et TD, D\'epartement Techniques de Commercialisation, 2\ieme{}~ann\'ee) \\
			& \textbf{Economie poursuite d'\'etudes} (CM et TD, D\'epartement Techniques de Commercialisation, 2\ieme{}~ann\'ee) \\
			& \textbf{A.T.A. Banque} (TD, D\'epartement Techniques de Commercialisation, 2\ieme{}~ann\'ee) \\
			& \textbf{Economie g\'en\'erale} (CM et TD, D\'epartement Techniques de Commercialisation, 1\iere{}~ann\'ee)
\end{longtable}

%%% CTU %%%

\smalltitre{Centre de T\'el\'e-Enseignement}

\noindent \begin{longtable}{@{}p{\annee}p{\texte}@{}}
2009 - 2015	& \textbf{Economie du d\'eveloppement} (CM, Master G\'eographie et Am\'enagement, sp\'ecialit\'e Am\'enagement et Gouvernance dans les Pays des Suds, 1\iere{}~ann\'ee) \\
\end{longtable}

\titre{U.F.R. de Sciences Economiques et de Gestion, Universit\'e Bourgogne Europe, Dijon}

\noindent \begin{longtable}{@{}p{\annee}p{\texte}@{}}
2020 -	2024 & \textbf{Quantitative evaluation of public policies} (CM, Master Data Analyst for Spatial and Environmental Economics, 2\ieme{}~ann\'ee) \\
2019 -	2023 & \textbf{Advanced topics in spatial statistics} (CM, Master Data Analyst for Spatial and Environmental Economics, 1\iere{}~ann\'ee) \\
2017 - 2018	& \textbf{Evaluation quantitative des politiques publiques} (CM, Master Economie Appliqu\'ee, 2\ieme{}~ann\'ee) \\
2016		& \textbf{Techniques d'enqu\^ete} (CM, Master Economie et gouvernance des territoires, 2\ieme{}~ann\'ee) \\
2015 - 2016	& \textbf{Choix publics et gouvernance des territoires} (CM, Master Economie et gouvernance des territoires, 2\ieme{}~ann\'ee) \\
			& \textbf{Traitement et analyse des donn\'ees localis\'ees} (CM, Master Territoires Environnement Energie, 1\iere{}~ann\'ee) \\
2001 - 2002 & \textbf{Dynamique macro\'economique} (TD, Ma\^itrise de Sciences Economiques et Ma\'etrise d'Econom\'etrie) \\
2000 - 2002 & \textbf{Statistiques descriptives} (TD, DEUG Economie et Gestion, 1\iere{}~ann\'ee) \\
1999 - 2000 & \textbf{Micro\'economie} (TD, DEUG Economie et Gestion, 2\ieme{}~ann\'ee) \\
1998 - 1999 & \textbf{Math\'ematiques} (TD, DEUG Economie et Gestion, 1\iere{}~ann\'ee) \\
1998 - 2002 & \textbf{Econom\'etrie} (TD, Licence de Sciences Economiques)
\end{longtable}

%%% Universit\'e Toulouse 1 %%%

\titre{Universit\'e Toulouse 1}

\noindent \begin{longtable}{@{}p{\annee}p{\texte}@{}}
2015 - 2019 & \textbf{Econom\'etrie spatiale} (CM, formation \`a distance, Master Economie et Statistique, sp\'ecialit\'e Statistique et Econom\'etrie, 2\ieme{}~ann\'ee) \\
\end{longtable}

%%% Universit\'e Montesquieu-Bordeaux IV %%%

\titre{Universit\'e Montesquieu-Bordeaux IV}

\noindent \begin{longtable}{@{}p{\annee}p{\texte}@{}}
2005 - 2006	& \textbf{Traitement des donn\'ees} (CM, Master Sciences Economiques et Sociales, mention Economique Appliqu\'ee : Histoire Economique, sp\'ecialit\'e recherche, 2\ieme{}~ann\'ee) \\
			& \textbf{Macro\'econom\'etrie} (CM, Master sp\'ecialit\'e Economiques et Sociales, mention Ing\'enierie Economique, sp\'ecialit\'e recherche, 2\ieme{}~ann\'ee) \\
2004 - 2006	& \textbf{Econom\'etrie et Econom\'etrie des donn\'ees de panel} (CM, Formation permanente de l'Ecole Doctorale en Sciences Economiques, Gestion et D\'emographie) \\
			& \textbf{Econom\'etrie II} (CM, Master Sciences Economiques et Sociales, mention Ing\'enierie Economique, 1\iere{}~ann\'ee, cours commun avec M.-A. S\'en\'egas) \\
			& \textbf{Econom\'etrie I} (CM, Master Sciences Economiques et Sociales, mention Ing\'enierie Economique, 1\iere{}~ann\'ee) \\
			& \textbf{Econom\'etrie} (CM, Magist\`ere d'Economie et de Finance Internationale, 2\ieme{}~ann\'ee) \\
			& \textbf{Introduction \`a l'\'econom\'etrie} (CM, Licence de Sciences Economiques, 3\ieme{}~ann\'ee, cours commun avec M.-A. S\'en\'egas) \\
2004 - 2005	& \textbf{Micro\'economie} (TD, Licence de Sciences Economiques, 2\ieme{}~ann\'ee) \\
2003 - 2006	& \textbf{Econom\'etrie spatiale} (CM, DEA Economie de l'Environnement, Innovation et Am\'enagement ; Master Sciences Economiques et Sociales, mention Economie Appliqu\'ee : Territoire, Environnement, Industries, sp\'ecialit\'e recherche) \\
			& \textbf{Statistiques} (CM, Magist\`{e}re d'Economie et de Finance Internationale, 1\iere{}~ann\'ee) \\
2003 - 2004	& \textbf{Econom\'etrie appliqu\'ee} (CM, Ma\^itrise Econom\'etrie) \\
			& \textbf{Statistiques} (CM, Ma\'etrise Econom\'etrie) \\
			& \textbf{Econom\'etrie appliqu\'ee} (CM, Licence d'Econom\'etrie) \\
			& \textbf{Econom\'etrie} (TD, Ma\^itrise de Sciences Economiques)
\end{longtable}

%%% Universit\'e d'Evry Val-d'Essonne %%%

\titre{Universit\'e d'Evry Val-d'Essonne}

\smalltitre{U.F.R. de Sciences Economiques et Juridiques}

\noindent \begin{longtable}{@{}p{\annee}p{\texte}@{}}
2003 - 2007 & \textbf{Econom\'etrie spatiale} (CM, M2 Ing\'enierie Economique et Statistique) \\
\end{longtable}

%%% Autres formations %%%

\titre{Autres formations}

\smalltitre{6\textsuperscript{th} Seminar Jean Paelinck}

\noindent \begin{longtable}{@{}p{\annee}p{\texte}@{}}
2013 & \textbf{Cours}, \og{}Using the R package for spatial econometrics\fg{} (19 octobre) \\
\end{longtable}

\smalltitre{CEPII}

\noindent \begin{longtable}{@{}p{\annee}p{\texte}@{}}
2013 & \textbf{Formation continue du CEPII}, \og{}Econom\'etrie spatiale\fg{} (11 juin) \\
\end{longtable}

\smalltitre{Toulouse School of Economics}

\noindent \begin{longtable}{@{}p{\annee}p{\texte}@{}}
2010 & \textbf{Spatial Econometrics School}, \og{}IV and GMM in spatial econometric models\fg{}, \og{}Spatial panels\fg{} (8-9 mars) \\
\end{longtable}

\end{document}

